\documentclass[acmsmall,screen]{acmart}

\newif \ifDraft \Draftfalse

\usepackage[scaled=0.78]{DejaVuSansMono} % use a better tt font for code
\usepackage[T1]{fontenc}

\usepackage{listings}
\usepackage{mathpartir}
\usepackage{mathtools}
%\usepackage{stmaryrd}
%\usepackage{pifont}
\usepackage{xspace}
\usepackage{util}
\usepackage{array}
\usepackage[inline]{enumitem}
\usepackage{xcolor}
\usepackage{tikz}
\usetikzlibrary{cd}
\usetikzlibrary{backgrounds}

%% Journal information
%% Supplied to authors by publisher for camera-ready submission;
%% use defaults for review submission.

\setcopyright{rightsretained}
\acmPrice{}
\acmDOI{10.1145/3571240}
\acmYear{2023}
\copyrightyear{2023}
\acmSubmissionID{popl23main-p325-p}
\acmJournal{PACMPL}
\acmVolume{7}
\acmNumber{POPL}
\acmArticle{47}
\acmMonth{1}
\received{2022-07-07}
\received[accepted]{2022-11-07}

\citestyle{acmauthoryear}

\lstset{
  basicstyle=\ttfamily,
  commentstyle=\color{darkgray}\itshape,
}

\lstdefinelanguage{Cogent}{
  morecomment=[l]{--},
  keywords={type, layout},
  keywordstyle=\bfseries,
  keywords=[2]{let, in, record, variant, at, after, using},
  keywordstyle=[2]\sffamily,
  %keywords=[3]{U8,U16,U32,U64,Bool,Unit,Char},
  %keywordstyle=[3]\color{teal}\bfseries,
  identifierstyle=\color{black},
  extendedchars=true,
  mathescape,
  breaklines=true,
  literate=
    {lambda}{$\lambda$}1
    {<<->>}{$\hole{-}\;$}1
    {->}{$\rightarrow$}1
    %{->}{$\rightarrow$}1
    %{<-}{$\leftarrow$}1
    %{|->}{$\mapsto$}1
    %{=>}{$\Rightarrow$}1
}

\lstdefinelanguage{isa}{
  columns=fullflexible,
  keepspaces,
  mathescape=true,
  morecomment=[s]{(*}{*)},
  morekeywords={locale,fixes,assumes,shows,and,lemma,definition,in,
    type_synonym,where,procedures,theorem,if,then,else,do,od,case,of,let},
  literate=
    {"}{}0
    {ÔøΩ}{$^\prime$}1
    {'}{$^\prime$}1
    {\\<^sub>6}{$_6$}1      
    {\\<^sub>4}{$_4$}1      
    {\\<^sub>3}{$_3$}1      
    {\\<^sub>2}{$_2$}1      
    {\\<^sub>1}{$_1$}1      
    {\\<^sub>0}{$_0$}1      
    {\\<^sub>f}{$_f$}1
    {\\<^sub>G}{$_G$}1      
    {\\<^sub>T}{$_T$}1
    {\\<forall>}{$\forall$}1
    {\\<exists>}{$\exists$}1
    {\\<equiv>}{$\equiv$}1
    {\\<Longrightarrow>}{$\Longrightarrow$}2
    {\\<longrightarrow>}{$\longrightarrow$}2
    {\\<Rightarrow>}{$\Rightarrow$}1
    {\\<rightarrow>}{$\rightarrow$}1
    {\\<leftarrow>}{$\leftarrow$}1
    {\\<longleftrightarrow>}{$\longleftrightarrow$}2
    {\\<Down>}{$\Downarrow$}2
    {\\<guillemotleft>}{$\mathopen{\{\mkern-4mu|}$}2
    {\\<guillemotright>}{$\mathclose{|\mkern-4mu\}}$}2
    {\\<And>}{$\bigwedge$}1
    {\\<and>}{$\land$}1
    {\\<or>}{$\lor$}1
    {\\<not>}{$\lnot$}1
    {\\<in>}{$\in$}1
    {\\<notin>}{$\notin$}1
    {\\<noteq>}{$\neq$}1
    {\\<le>}{$\le$}1
    {\\<ge>}{$\ge$}1
    {\\<cup>}{$\cup$}1
    {\\<cap>}{$\cap$}1
    {\\<lambda>}{$\lambda$}1
    {\\<times>}{$\times$}1
    {\\<turnstile>}{$\vdash$}1
    {\\<turnstile-t>}{$\vdash_t$}2
    {\\<subseteq>}{$\subseteq$}1
    {\\<sqinter>}{$\sqcap$}1
    {\\<lbrace>}{$\mathopen{\{\mkern-4mu|}$}1
    {\\<rbrace>}{$\mathclose{|\mkern-4mu\}}$}1
    {\\<lbrakk>}{$\mathopen{\lbrack\mkern-3mu\lbrack}$}1
    {\\<rbrakk>}{$\mathclose{\rbrack\mkern-3mu\rbrack}$}1
    {\\<infinity>}{$\infty$}1
    {\\<sigma>}{$\sigma$}1
    {\\<mu>}{$\mu$}1
    {\\<gamma>}{$\gamma$}1
    {\\<tau>}{$\tau$}1
    {\\<Gamma>}{$\Gamma$}1
    {\\<le>}{$\leq$}1
    {\\<Sum>}{$\sum$}1
    {\\<dollar>}{\$}1
    {\\<lparr>}{$\mathopen{(\mkern-4mu|}$}1
    {\\<rparr>}{$\mathclose{|\mkern-4mu)}$}1
    }



\renewcommand{\sectionautorefname}{Section}
\renewcommand{\subsectionautorefname}{Section}
\newcommand\litem[1]{\item{\bfseries [#1]\enspace}}

\newcommand{\citepos}[2][]{\citeauthor{{#2}}{#1}~\cite{{#2}}}

% Shorthands
\newcommand{\cc}{\lstinline[language=C]}
\newcommand{\cg}{\lstinline[language=Cogent]}
\newcommand{\key}[1]{\texttt{{#1}}}
\newcommand{\bd}[1]{\textbf{{#1}}}
\newcommand{\fd}[1][f]{\key{{#1}}} % field
\newcommand\GG{\Gamma}
\newcommand\LL{\mathcal{L}}
\newcommand{\Sg}{\Sigma}
\newcommand{\sg}{\sigma}
\newcommand\adj{\delta}
\newcommand{\Adj}{\Delta}
\newcommand\bN{\mathbb{N}}
\newcommand\Ge{\varepsilon}
\newcommand\Go{\omega}
\newcommand\Gt{\theta}
\newcommand\Gg{\gamma}
\newcommand\GL{\Lambda}
\newcommand\Gl{\lambda}
\newcommand\rB{\mathrm{B}} % Byte descriptor
\newcommand\rb{\mathrm{b}} % bit descriptor
\newcommand{\fresh}[1]{{{#1}~\text{fresh}}}
\newcommand{\sem}[1]{\llbracket #1 \rrbracket }

\newcommand{\Cogent}{\textsc{Cogent}\xspace}
\newcommand{\cogent}{\Cogent}
\newcommand{\Dargent}{\textsc{Dargent}\xspace}
\newcommand{\dargent}{\Dargent}
\newcommand{\bilby}{BilbyFs\xspace}


\ifDraft
  \newcommand{\Comment}[1]{\textbf{\textsl{#1}}}
  \newcommand{\TODO}[1]{\textbf{\textsl{\colorbox{yellow}{TODO:} #1}}}
  \newcommand{\FIXME}[1]{\textbf{\textsl{\colorbox{yellow}{FIXME:} #1}}}
  \usepackage{draftwatermark}
  \SetWatermarkLightness{0.85}
\else
  \newcommand{\Comment}[1]{\relax}
  \newcommand{\TODO}[1]{\relax}
  \newcommand{\FIXME}[1]{\relax}
\fi
\newcommand{\zilinc}[1]{\Comment{#1 [zilinc]}}
\newcommand{\liam}[1]{\Comment{#1 [liam]}}
\newcommand{\ambroise}[1]{\Comment{#1 [ambroise]}}
\newcommand{\louis}[1]{\Comment{#1 [louis]}}
\newcommand{\cmcl}[1]{\Comment{#1 [cmcl]}}
\newcommand{\review}[1]{\Comment{#1 [reviewer]}}


\newcommand{\urltimerdriver}{https://github.com/seL4/util_libs/blob/c446df1f1a3e6aa1418a64a8f4db1ec615eae3c4/libplatsupport/src/plat/odroidc2/meson_timer.c}
\newcommand{\urltimerdriverh}{https://github.com/seL4/util_libs/blob/c446df1f1a3e6aa1418a64a8f4db1ec615eae3c4/libplatsupport/plat_include/odroidc2/platsupport/plat/meson_timer.h}
\hyphenation{poly-morphic}
\hyphenation{data-type}

\begin{document}
\title{\Dargent: A Silver Bullet for Verified Data Layout Refinement}


%% Author with single affiliation.
\author{Zilin Chen}
% \authornote{with author1 note}          %% \authornote is optional;
                                        %% can be repeated if necessary
\orcid{0000-0003-0854-2464}             %% \orcid is optional
\affiliation{
  %\position{Position1}
  %\department{Department1}              %% \department is recommended
  \institution{UNSW Sydney}            %% \institution is required
  %\streetaddress{Street1 Address1}
  %\city{City1}
  %\state{State1}
  %\postcode{Post-Code1}
  \country{Australia}                    %% \country is recommended
}
\email{zilin.chen@student.unsw.edu.au}          %% \email is recommended

%% Author with two affiliations and emails.
\author{Ambroise Lafont}
%\authornote{with author2 note}          %% \authornote is optional;
                                        %% can be repeated if necessary
\orcid{0000-0002-9299-641X}             %% \orcid is optional
\affiliation{
  %\position{Position2a}
  %\department{Department2a}             %% \department is recommended
  \institution{University of Cambridge}           %% \institution is required
  %\streetaddress{Street2a Address2a}
  \country{UK}                   %% \country is recommended
}
\email{ael62@cam.ac.uk}         %% \email is recommended

\author{Liam O'Connor}
\orcid{0000-0003-2765-4269}
\affiliation{
  \institution{University of Edinburgh}
  \country{UK}
}
\email{l.oconnor@ed.ac.uk}

\author{Gabriele Keller}
\orcid{0000-0003-1442-5387}
\affiliation{
  \institution{Utrecht University}
  \country{Netherlands}
}
\email{g.k.keller@uu.nl}

\author{Craig McLaughlin}
\orcid{0000-0002-1323-8566}
\affiliation{
  \institution{UNSW Sydney}
  \country{Australia}
}
\email{c.mclaughlin@unsw.edu.au}

\author{Vincent Jackson}
\orcid{0000-0002-8737-4202}
\affiliation{
  \institution{University of Melbourne}
  \country{Australia}
}
\email{vjjac@student.unimelb.edu.au}

\author{Christine Rizkallah}
\orcid{0000-0003-4785-2836}
\affiliation{
  \institution{University of Melbourne}
  \country{Australia}
}
\email{christine.rizkallah@unimelb.edu.au}



%% Abstract
%% Note: \begin{abstract}...\end{abstract} environment must come
%% before \maketitle command

\begin{abstract}
%Systems programming involves writing code which adheres to well-defined
%interfaces usually imposed by standardised protocols or hardware. 
%\cogent's certifying compiler generates imperative C programs and
%provides a formal proof that its compilation is a refinement of the
%functional \cogent code. 

Systems programmers need fine-grained control over the memory
layout of data structures, both to produce performant code and to comply with well-defined interfaces imposed by existing code, standardised protocols or hardware. Code that manipulates these low-level representations in memory is hard to get right. Traditionally, this problem is addressed by the implementation of tedious marshalling code to convert between compiler-selected data representations and the desired compact data formats. Such marshalling code is error-prone and can lead to a significant runtime overhead due to
excessive copying. While there are many languages and systems that address the correctness issue, by automating the generation and, in some cases, the verification of the marshalling code, the performance overhead introduced by the marshalling code remains.
In particular for systems code, this overhead can be prohibitive.
In this work, we address both the correctness and the performance problems.

We present a data layout description language and data refinement framework,
called \Dargent, which allows programmers to declaratively specify
how algebraic data types are laid out in memory.
Our solution is applied to the \Cogent language,
but the general ideas behind our solution are applicable to other settings.
The \Dargent framework generates C code that manipulates data directly with the desired memory layout,
while retaining the formal proof that this generated C code is correct
with respect to the \cogent functional semantics.
This added expressivity removes the need for implementing and verifying marshalling code,
which eliminates copying, smoothens interoperability with surrounding systems,
and increases the trustworthiness of the overall system.
\end{abstract}



\begin{CCSXML}
<ccs2012>
   <concept>
       <concept_id>10011007.10011074.10011099.10011692</concept_id>
       <concept_desc>Software and its engineering~Formal software verification</concept_desc>
       <concept_significance>500</concept_significance>
       </concept>
   <concept>
       <concept_id>10011007.10011006.10011060.10011690</concept_id>
       <concept_desc>Software and its engineering~Specification languages</concept_desc>
       <concept_significance>500</concept_significance>
       </concept>
   <concept>
       <concept_id>10011007.10010940.10010941.10010949</concept_id>
       <concept_desc>Software and its engineering~Operating systems</concept_desc>
       <concept_significance>500</concept_significance>
       </concept>
   <concept>
       <concept_id>10011007.10011006.10011008.10011024.10011028</concept_id>
       <concept_desc>Software and its engineering~Data types and structures</concept_desc>
       <concept_significance>300</concept_significance>
       </concept>
 </ccs2012>
\end{CCSXML}

\ccsdesc[500]{Software and its engineering~Formal software verification}
\ccsdesc[500]{Software and its engineering~Specification languages}
\ccsdesc[500]{Software and its engineering~Operating systems}
\ccsdesc[300]{Software and its engineering~Data types and structures}


%%
%% Keywords. The author(s) should pick words that accurately describe
%% the work being presented. Separate the keywords with commas.
\keywords{certifying compiler, data refinement, systems programming}

\maketitle

\renewcommand{\shortauthors}{Z. Chen, A. Lafont, L. O'Connor, G. Keller, C. McLaughlin, V. Jackson, and C. Rizkallah}

% \balance


\section{Introduction}\label{sec:intro}

In the realm of software systems, such as device drivers, file
systems, and network stacks, precise control over the
\emph{data layout} of objects is crucial for compatibility and
performance. Specifically, controlling the composition of objects in memory
on a bit- and byte-level can avoid the need for translation or \emph{deserialisation} 
at the boundaries between on-medium and
in-memory data which frequently arise from interacting with
standardised protocols or hardware.
%
These systems are often implemented in the C language, in part because it offers low-level features to give fine-grained control over data layout. 
Unfortunately, to maintain good performance, the C programmer must throw away the conceptual abstraction of the data type, and instead focus on the low-level details of bits and bytes. This low-level code contains many subtle bit-twiddling
operations which, apart from being difficult to manually verify, are
also tedious and error-prone to implement.

% 
In Linux, a \cc{dev_t} type is used to represent device numbers. It is \cc{typedef}ed to
a 32-bit unsigned integer (\cc{u32}). It is composed of two parts: a major number and a minor
number. Within the 32 bits, 12 bits are set aside for the major number, and the remaining
20 bits for the minor number. In C, to construct a \cc{dev_t} or to fetch the device
numbers, users need to use a some macros, which are implemented with convoluted bit twiddling:
\begin{lstlisting}[language=C]
#define MINORBITS    20
#define MINORMASK    ((1U << MINORBITS) - 1)
#define MAJOR(dev)   ((unsigned int) ((dev) >> MINORBITS)) // get major
#define MINOR(dev)   ((unsigned int) ((dev) & MINORMASK))  // get minor
#define MKDEV(ma,mi) (((ma) << MINORBITS) | (mi))          // construct dev_t
\end{lstlisting}
Logically, the \cc{dev_t} type is simply a product of two integers.
To implement it in a high-level language, users would typically define it as a record
type with two integer fields. For example, in \cogent, it may be defined as:\footnote{
As an opening example, we show pseudo-\cogent code here to omit the irrelevant \cogent
intricacies.}

\begin{lstlisting}[language=Cogent]
type Dev_t = { major : Int, minor : Int }
\end{lstlisting}
and to access the fields of such an object \cg{dev},
the normal member operation (\cg{dev.major} and \cg{dev.minor}) can be used.
With this abtract view of the type, which facilitates simpler operations on this type,
the user unfortunately has to hand over control of the type's layout to the
\cogent compiler. The user then needs to study the compiler's code generation convention
and write data format conversion functions to related the compiled \cg{Dev_t} type with
the native C type \cc{dev_t}. In fact, the user does not have to make this compromise. With
\dargent, the user can have both the abstraction and control over types' memory
layouts. By specifying a \dargent layout description, the \cogent compiler can
generate C code with the specified layout:
\begin{lstlisting}[language=Cogent]
type Dev_t = { major : Int, minor : Int }
  layout record { major : 12b, minor : 20b }
\end{lstlisting}
This is transparent to the functional smenatics
of the \cogent program. Namely, the \cg{dev.major} operation continues to work regardless
of the underlying layout of \cg{dev}.





To maintain the higher-level structure of a program without sacrificing
performance, we want to use a language with a high-level semantics, but with facilities for
specifying the low-level memory layout of heap-allocated objects. While most
high-level languages use fixed heap layouts, there has been recent
progress on language-based support for user-defined memory
layouts~\cite{Cronburg:2019,Vollmer:2019}. These languages, however, do
not provide verified correctness of their translations.

In this paper, we present \Dargent, a language for describing data
layouts of high-level algebraic datatypes along with a data
refinement framework for \emph{automatically verifying} the correctness of the compiled C code
with respect to those layout descriptions. We
build on the \cogent language and refinement framework~\citep{Cogent_JFP_2021}.
\cogent is designed for the implementation of high-assurance low-level systems components as pure
mathematical functions operating on algebraic data types.
\cogent's certifying compiler co-generates
a C program and Isabelle/HOL~\citep{Nipkow:2002} theorems witnessing a proof
that the C program refines an Isabelle/HOL embedding of the \cogent source program~\cite{Rizkallah:itp16}. 

While there is a very long line of prior work on \emph{data description
languages}~\citep{FisherG05, McCannC00, Back02, WangG11, everparse,
Ye_Delaware_19, vanGeest_Swierstra_17, Slind_21}, and the data layout
descriptions used in \dargent do indeed look similar to those used in such
languages, there is a fundamental difference: These languages are
designed for synthesising data (de)serialisation functions (also sometimes
referred to as data marshalling/unmarshalling functions, encoders/decoders, 
or parsers and pretty-printers for low-level data),
which convert data stored in a low-level, sequential format in some storage
medium to a high-level, structured representation in memory and vice versa. 
They are primarily used for the interaction and communication between
different programming languages (e.g.\ a foreign function interface)
or systems (e.g.\ data transmission over the network). In this context, code to transform back-and-forth
between the two representations is still necessary.

\dargent, on the other hand, is intended to solve a different
problem. The \dargent data layout descriptions grant programmers
the ability to dictate to the \cogent compiler how it should lay out the algebraic data
types used by the \cogent program itself. The compiler generates code that works directly 
with data laid out according to the programmer's specifications, as well as Isabelle/HOL  
proofs showing that it has done so correctly. 

Depending on the application, \dargent therefore can either eliminate entirely
or reduce the need for data (de)serialisation code, be it manually
written or automatically derived, when interacting with the external world.
Because algebraic data types can be represented directly in their binary data formats with \dargent,  the 
programmer does not need
to decode raw data first into some other in-memory representation in order to operate on it as a data type. Eliminating these (de)serialisation steps naturally results in more concise and readable code, better performance, and easier informal and 
formal reasoning.

This additional power in expressiveness can also be
used to improve the performance of the compiled \cogent code, 
e.g.\ by having smaller memory footprints or a specialised memory
layout optimal for the underlying architecture,
independent of the interoperation between languages or systems.
It also enables \cogent programmers to directly write code that is binary-compatible
with native C programs.

\cogent is readily amenable to an extension for prescribing 
data layouts and fine-tuning the compilation of algebraic data types 
by virtue of its lack of a runtime system
and direct compilation to C. Before \dargent, it simply adopted
the layout conventions of the
underlying C compiler. The introduction of \dargent into the
framework enabled improvements to
some outstanding inefficiencies in the prior design, such as
a reduced reliance on deserialisation code within file system
implementations~\citep{Amani:2016} and directly representing device
register formats as data types.
%%
Furthermore, we have extended the \cogent compiler so as to preserve the
benefits of \cogent's high-level type system and semantics. The upshot
is our compiler automatically translates read and write operations on
heap-allocated objects to take account of their particular data
layout. The translation's correctness is guaranteed by the enhanced
data refinement theorem. 
This paper describes, to the best of our knowledge, the first framework that
is able to leverage data layout specifications for generating bit-level accessors
with formal correctness guarantees. %Thanks to \cogent's uniqueness type system, our framework also supports mutable data. 
%This is the first \emph{certifying} framework to be extended with data layout descriptions supporting mutable data, thanks to \cogent's
%uniqueness type system which allows for efficient destructive
%updates.

\subsection*{Contributions}
This paper realises the vision set out previously~\citep{Isola:2018}.
We make the following contributions:
\begin{itemize}
  \item The design and implementation of \Dargent
    (\autoref{sec:dargent}), a \emph{data layout} language for
    controlling the memory layout of algebraic data types, down to the
    bit level. We formalise a core calculus and its static semantics
    in Isabelle/HOL~\citep{Nipkow:2002}, and discuss the compilation process;
  \item An extension to the \cogent verification framework 
    (\autoref{sec:verification}) to \emph{automatically verify} the
    translation of high-level read/write accesses (known as
    \emph{getters} and \emph{setters}) to explicit offsets within a
    well-defined memory region;
  \item An extended suite of examples (\autoref{sec:applications}
    and \autoref{sec:timerdriver}) demonstrating the utility of \Dargent in
    the context of device drivers.
\end{itemize}
All the results described in this paper, including the case studies we
present, are associated with formal proofs in Isabelle/HOL.
These materials are packaged as a virtual machine image, which is publicly
available~\citep{dargent-artefact}. It is derived from a development branch of
the \cogent project repository (\url{https://github.com/au-ts/cogent}).


\section{An Overview of \Cogent}{\label{sec:cogent}}

\Cogent~\citep{OConnor:2016, OConnor:thesis, Cogent_JFP_2021} is a higher-order, polymorphic, purely functional programming
language, in the tradition of the pure subsets of languages such as
Haskell or ML. Programs are expressed as mathematical functions
operating on algebraic data types. Unlike Haskell or ML, however,
\Cogent is designed for implementing low-level systems code where manual
memory management is essential for performance reasons. As such,
\Cogent has no garbage collection mechanism and heap (de)allocations are
explicit in the language.

The \cogent language is equipped with a \emph{uniqueness type
system}~\citep{wadler:90, clean:93}, enabling a seamless transition from a
purely functional semantics to an imperative semantics and a compilation
strategy to generate efficient low-level C code.
%
The certifying compiler automatically produces a formal
proof~\citep{Rizkallah:itp16} that its generated C code refines an Isabelle/HOL
embedding of the \cogent program's functional semantics. \Cogent was used to
implement two real-world file systems and to verify the correctness of key file
system operations~\citep{Amani:2016}.
%The \cogent language and its verification were introduced
%by~\citepos{OConnor:2016,} and most recently presented by
%\citeauthor{OConnor:thesis} in his PhD
%thesis~\cite{OConnor:thesis}. Real-world applications using \cogent for
%certifying file systems are presented by~\citepos{Amani:2016} (see also
%\citepos['s PhD thesis]{Amani:phd}). 
This section summarises aspects of \cogent and its verification framework that
are relevant to our work.

\subsection{\Cogent's Uniqueness Types System}

Uniqueness type systems ensure that each \emph{linear} object in memory
is uniquely referenced. Consequently, updates to these objects in a
purely functional language can be compiled as in-place destructive
updates, without the need for copying~\citep{wadler:90}. In \cogent, we
call the type of objects that are subject to the uniqueness restrictions
\emph{linear types}, and the rest \emph{non-linear types}.  Roughly
speaking, any object that resides in the heap or contains pointers to
other heap-objects --- and is not \emph{read-only} (explained later) ---
is linear; otherwise it is non-linear.  In the generated C code, all
pointers address heap memory.  This is because \cogent's verification
framework depends on the AutoCorres library~\citep{autocorres}, which
does not support stack pointers.  Therefore, we can summarise the
linearity of an object simply as: A linear object is behind a pointer
and/or contains pointers.
%% Later we will see that the linearity of a type is closely correlated
%% to the \dargent layouts.

As a simple \cogent example, consider the following code.
\begin{lstlisting}[language=Cogent]
  type Bag = { count : U32, sum : U32 }
  
  addToBag : (U32, Bag) -> Bag
  addToBag (x, bag {count = c, sum = s}) = bag {count = c + 1 , sum = s + x }
\end{lstlisting}
The first line declares a linear record type \cg|Bag|, which is a record
comprised of two 32-bit unsigned words. In the \cg|addToBag| function,
we retrieve the value of the two fields by pattern matching and 
update both fields of the record according to the given input.
It is compiled to a C function that takes as input a pointer to a bag
structure and updates it \emph{in place}, because the uniqueness type 
system ensures that the \cg|Bag| object is uniquely referenced.


\begin{figure}
\[
\begin{array}{llclr}
  \text{Types} 
     & \tau & ::=  & T & \text{(primitive types)} \\
     &      & \mid & ()                & \text{(unit)} \\
     &      & \mid & \alpha            & \text{(type variables)} \\
     % &      & \mid & \tau_1 \rightarrow \tau_2      & \text{(functions)} \\
     &      & \mid &\RecordTy{\many{\mathrm{f}_i : \tau_i}}{s} & \text{(records)}\\
     &      & \mid & \VariantTy{\many{\mathrm{A}_i\ \tau_i}}  & \text{{(variants)}}\\
     &      & \mid & \AbsTy{K}{\many{{\tau_i}}}{s}  & \text{(abstract types)} \\
     &      & \mid & \dots & \\
  \text{Primitive types} & T & ::= & \PrimType{U}n \mid \PrimType{Bool} & \text{($n = 1 \dots 64$)} \\
  \text{Field names}  & & \ni  & \mathrm{f}\\
  \text{Constructors} & & \ni  & \mathrm{A}\\
  \text{Sigils} & s & ::= & \BoxedNL \mid \Unboxed & \\
                & \BoxedNL & ::= & \ModeN{w} \mid \ModeN{r} & 
\end{array}
\]
(lists are represented by $\many{\text{overlines}}$)
\caption{The syntax of \cogent types}
\label{fig:syntax-cogent-types}
\end{figure}


The syntax of (a relevant subset of) \cogent's types is given in \autoref{fig:syntax-cogent-types}.
\cogent \emph{primitive} types include the $n$-bit unsigned integer
types, $\PrimType{U}n$ (where $n = 1 \dots 64$), and booleans ($\PrimType{Bool}$).
Among the $\PrimType{U}n$ types, we call them \emph{word-sized integers} (or simply
\emph{words}) when $n \in \{8,16,32,64\}$. Algebraic data types can be formed using
\emph{record} and \emph{variant} types (aka.\ tagged unions).
%
Furthermore, \cogent supports declaring \emph{abstract types} with their
definitions provided in C.
%
\cogent employs a structural type system, meaning that compound types, as well as their subtyping
relations, are defined solely by their structure and the structure of their components. Names given to types such as $\mathtt{Bag}$ above are type synonyms---mere abbreviations for better cosmetics. 


\cogent's type system distinguishes between \emph{boxed} and
\emph{unboxed} types through a \emph{sigil} annotation on the type.
The boxedness of a type characterises how an object is referenced---by-value or by-pointer.
An object of a boxed type (sigil $\BoxedNL$) is accessed via its pointer.
Because in \cogent, all pointers address heap memory, a boxed object must reside in the heap.
An object of an unboxed type (sigil $\Unboxed$) is accessed by value and either
resides in the stack or is embedded inside a larger data structure in the
heap. In the latter case, when the object is accessed, it will be copied
by-value to a stack variable. The sigil is only used in the \cogent core calculus.
For example, the type \cg{Bag} we showed earlier is desugared to $\RecordTy{
  \FieldTy{count}{\PrimType{U32}}, \FieldTy{sum}{\PrimType{U32}}}{\BoxedNL}$, as it is 
a boxed record.

The uniqueness and the boxedness of a type are concerned with different
aspects of memory locations.  The uniqueness relates to \emph{what} kind
of memory the object references (linear or non-linear), whereas the
boxedness of an object determines \emph{how} its memory is referenced
(by-value or by-pointer).
%
Since all pointers in \cogent point to the heap, any boxed object must
be linear (with the same caveat that it must not be read-only) because
it is directly behind a pointer. The converse is not true though. A
typical counter-example is a stack-allocated structure that contains
pointers; this structure is linear but unboxed.

Records and abstract types can be either boxed or unboxed. Primitive types and
variant types, however, can only be unboxed. For this reason, primitives and
variants do not include a sigil. For a boxed type, \cogent allows further
fine-grained control of accessibility: A boxed type can either be writable
(with sigil $\ModeN{w}$) or read-only (with sigil $\ModeN{r}$). When a type
has a read-only sigil, even though it is behind a pointer or contains pointers,
it becomes non-linear---this is the caveat we mentioned earlier in this section.
This read-only object can be shared analogously to how a variable can be immutably borrowed
in the type system of Rust~\citep[\S~4.2]{rust-borrow}.

Before \dargent{}, the \cogent compiler used a pre-defined code generation algorithm to compile
data types to C.
A record type in \cogent was mapped to a C
struct, with the fields laid out in the order in which they are declared in the
type. The mapping for variant types, on the other hand, was less direct:
\Cogent's verification tool chain does not support C unions, so variants were also represented as structs in C, containing a field for a tag, and a field for each alternative's payload.

Applying such a fixed code generation scheme is no surprise for a typical high-level
functional language,
whose implementation details are hidden from the language users. But \cogent is not just a functional language, it is also a systems language where
the exact low-level representation of data types is relevant to programmers.
This fixed code generation algorithm often resulted in
suboptimal or undesired representations of \cogent types in C, and users of the 
\cogent language had to know
about the implementation details of the code generator in order to write
C code that directly interfaces with \cogent programs. In situations where
the \cogent program needs to interoperate with existing C components, say,
when developing an operating system component that interfaces with the Linux
kernel headers, glue code was required to translate between representations, resulting in development
and run-time overhead. \liam{Perhaps we should add some Bilby numbers here}
%
%% CMCL: Reword sentence below
%%
Also, problematically, as this glue code depended on the representation choices
made by the \cogent compiler, future versions of the compiler could break
previously working code due to changes in the code generation scheme.

\subsection{The Verification Framework}
{\label{ssec:cogent-verification}}

In addition to a programming language, \Cogent is also a verification framework
realised in Isabelle/HOL, based on \emph{certifying compilation}.
Compiling a \cogent program results in multiple components:
\begin{enumerate}
  \item a C program,
  \item an Isabelle/HOL \emph{shallow embedding}   of the \cogent program,
  \item an Isabelle/HOL \emph{deep embedding}   of the \cogent program,
  \item an Isabelle proof of refinement between the C
      program and the Isabelle shallow embedding.
\end{enumerate}
The last item relies on the AutoCorres library~\citep{autocorres} to generate 
a representation of the C code in Isabelle/HOL. More precisely, the AutoCorres library 
abstracts the C semantics via a SIMPL~\citep{simpl} formal language into
a monadic embedding in Isabelle/HOL.

This compilation process provides an indirect way of formally verifying
properties about the generated C program. The user first proves
the desired properties about the Isabelle shallow
embedding.  This manual proof should follow by simple equational reasoning about HOL terms by applying term rewriting
tactics provided by the theorem prover. Then, the automatic refinement
proof between the C code and the shallow embedding transports the proven
properties to the C program. In short, \cogent's verification
framework reduces complicated low-level verification on the C program to
a simple high-level equational proof.

The compiler-generated refinement proof between the C code and the
Isabelle shallow embedding is composed of several smaller
correspondence proofs.  When the compiler generates the shallow
embedding of the \cogent program in Isabelle, it also generates a
\emph{deep embedding} representing the abstract syntax of the \cogent
program. Two semantics are assigned to the deep embedding: a purely
functional \emph{value semantics}, and a stateful \emph{update
  semantics}, with pointers, memory states, and in-place field updating.
It is proved once-and-for-all that these two semantics are equivalent for
any well-typed \cogent program. As part of the certifying compilation,
the compiler produces a refinement proof between the shallow embedding
and the deep embedding with value semantics, and a refinement proof
between the deep embedding with update semantics and the monadic C
embedding obtained from AutoCorres. Chaining the three correspondence
lemmas results in a correspondence between
the shallow embedding and the C program, stating that the C program is a
refinement of the \cogent program's semantics (see \autoref{fig:cogent_framework}).

\begin{figure}
\begin{tikzpicture}
  \draw (0,0.5) rectangle (3.5,1.5) node[pos=0.5, text width=3.5cm, align=center] 
        {\textbf{Shallow Embedding} \\ \small (pure functions)};
  \draw (-0.75,-0.7) rectangle (4.25,-3.1) node[xshift=0.4cm, pos=0.5, text width=3.5cm, align=center, anchor=east] 
        {\textbf{Deep\\Embedding}};
  \draw [fill=green!10] (0.5,-0.8) rectangle (3.3,-1.3) node[pos=0.5, text width=3.5cm, align=center] 
        {{\small Value Semantics}};
  \begin{scope}[yshift=2cm - 3.5pt]
    \draw  (1.35,-4.25) -- (1.35,-3.85) -- (1.2,-3.85) -- (1.75,-3.25) -- (2.3,-3.85) -- (2.15,-3.85)
           -- (2.15,-4.25) -- (1.35,-4.25);
    \node [text width=1cm, align=center] at (1.74,-3.7) {\large ${\sqsubseteq}$};
    \node [text width=1cm, align=center] at (1.75,-4) {\scriptsize refines};
  \end{scope}

  \draw [fill=red!10] (0.5,-2.5) rectangle (3.3,-3) node[pos=0.5, text width=3.5cm, align=center] 
        {{\small Update Semantics}};
  \draw (0,-5.3) rectangle (3.5,-4.3) node[pos=0.5, text width=3.5cm, align=center] 
        {\textbf{C Code} \\ \small (via AutoCorres)};
  \begin{scope}[xshift=0,yshift=5pt]
    \draw  (1.35,-0.75) -- (1.35,-0.35) -- (1.2,-0.35) -- (1.75,0.25) -- (2.3,-0.35) -- (2.15,-0.35)
           -- (2.15,-0.75) -- (1.35,-0.75);
    \node [text width=1cm, align=center] at (1.74,-0.2) {\large ${\sqsubseteq}$};
    \node [text width=1cm, align=center] at (1.75,-0.5) {\scriptsize refines};
  \end{scope}
  \begin{scope}[yshift=1.5pt]
    \draw  (1.35,-4.25) -- (1.35,-3.85) -- (1.2,-3.85) -- (1.75,-3.25) -- (2.3,-3.85) -- (2.15,-3.85)
           -- (2.15,-4.25) -- (1.35,-4.25);
    \node [text width=1cm, align=center] at (1.74,-3.7) {\large ${\sqsubseteq}$};
    \node [text width=1cm, align=center] at (1.75,-4) {\scriptsize refines};
  \end{scope}
\end{tikzpicture}

\caption{The \cogent refinement framework.}
\label{fig:cogent_framework}
\end{figure}


\section{\Dargent}{\label{sec:dargent}}

We have designed and implemented a data layout description language called \dargent,
which describes how a \cogent algebraic data type may be laid
out in (heap) memory, down to the bit level. Layout descriptions in \dargent are transparent to
the shallow embedding of \cogent's semantics, but they
influence the definition of the refinement relation to C code
generated by the compiler. In \autoref{sec:verification}, we describe in more
detail the extensions to our verification framework to accommodate
the \dargent layout descriptions. Here, we focus on the language
definition. Firstly, we give an informal overview of \dargent's language features.

\subsection{An Informal Introduction to \dargent} \label{ssec:informal-intro}

%%
%% CMCL: Reword; type vs. object distinction
%%
\Dargent offers the possibility to assign any boxed \cogent type a
custom layout describing how an object of that type should be stored in
the heap.
%
A boxed type
assigned with a custom layout is compiled to a C struct with a single
field: an array of 32-bit words.\footnote{Throughout the paper, unless
we explicitly specify the size, the term
``word'' always refers to unsigned integer types of 1 byte, 2 bytes,
4 bytes or 8 bytes and the actual size is usually less relevant in
the discussion. It does not necessarily imply pointer-sized words.
We discuss the implications
of the choice of the word size later in \autoref{ssec:gsetters}. 
}
This array \emph{represents} the
\cogent type: It can be deemed as untyped in C, but the \dargent description
contains enough information to access individual parts of the type correctly.
How data is laid out in memory is purely a low-level concern, and it does not
affect the functional semantics of a \cogent program in any way.
In other words, \cogent functions are parametric over the layouts of types.
%
Under the hood, the \cogent compiler generates custom getters (and
setters) in C, retrieving (and setting) the relevant parts of the object
from its array representation.

As an example, consider the \cogent type in \autoref{fig:cogent-type-ex}.
\begin{figure}
  \centering
  $$
  \begin{array}{ll}
    \mathbf{type}\ \AbsN{Example} = \{ \\
  \quad \FieldTy{struct}{\#\{\FieldTy{x}{\PrimType{U32}}, \FieldTy{y}{\PrimType{Bool}}\}},
    & \!\!\textit{-{}- nested embedded record} \\
  \quad \FieldTy{ptr}{\{\dots\}},
    & \!\!\textit{-{}- pointer to another record} \\
  \quad \FieldTy{sum}{\VariantTy{\Cons{X}{\PrimType{U16}} \alt \Cons{Y}{\PrimType{U8}}}}
    & \!\!\textit{-{}- variant type} \\
    \} 
  \end{array}
  $$
  \caption{\Cogent type example}
  \label{fig:cogent-type-ex}
\end{figure}
This record consists of three fields: $\FieldN{struct}$, $\FieldN{ptr}$ and $\FieldN{sum}$.
The type of the $\FieldN{struct}$ field is an unboxed record, denoted by the leading $\#$ symbol. This field is
embedded inside the parent $\AbsN{Example}$ record. The $\FieldN{ptr}$ field is a boxed
record, and is stored somewhere else in the heap, referenced by a pointer in the $\AbsN{Example}$
type. The last field $\FieldN{sum}$ has a variant type, with two alternatives tagged $\ConsN{A}$ and $\ConsN{B}$
respectively. This variant is unboxed (recall that there is no boxed variant in \cogent) and
is stored in the heap inside the parent record.

A layout for this record type must specify where each field is located in the
word array. Overlapping is not allowed, except for the payloads of the two constructors
of the variant type, since only one of them is
relevant at each time, depending on the tag value. The layout must also specify 
what the tag values for the variant constructors $\ConsN{A}$ and $\ConsN{B}$ are.
We will give a layout to this type after a short
introduction to the \dargent language constructs.

% Layouts cannot currently be assigned to \emph{unboxed} records, i.e., records that
% are not behind a pointer (for example, when declaring a record as a local
% variable). We plan to overcome this limitation in the future.
%

\begin{figure}
$$
  \begin{array}{llclr}
  \text{Sizes} & s & ::= & \multicolumn{2}{l}{n\rB \mid n\rb \mid s + s} \\
  \text{Layout expressions} 
     & \ell & ::=  & s & \text{(block of memory)}\\
     &      & \mid & x & \text{(layout variable)} \\
     &      & \mid & \mathsf{pointer} & \text{(pointer layout)}\\
     &      & \mid & L\ \many{\ell_i} & \text{(another layout)}\\
     &      & \mid & \ell\ \mathsf{at}\ s & \text{(offset operator)}\\
     &      & \mid & \ell\ \mathsf{after}\ \FieldN{f} & \text{(relative location)}\\
     &      & \mid & \ell\ \mathsf{using}\ \omega & \text{(endianness)}\\
     &      & \mid & \multicolumn{2}{l}{\mathsf{record}\ \{ \overline{\FieldN{f}_i : \ell_i} \}} \\
     &      & \mid & \multicolumn{2}{l}{\mathsf{variant}\ (\ell)\ \{ \overline{\ConsN{A}_i\ (n_i) : \ell_i} \}} \\
     % &      & \mid & \multicolumn{2}{l}{\mathsf{array}\ \{ \ell \}} \\
  \text{Declarations} & d & ::= & \multicolumn{2}{l}{\mathbf{layout}\ L\ \many{x_i} = \ell} \\
  \text{Layout names} &   & \ni & L\\
  \text{Endianness} & \omega & ::= & \bigEnd \mid \littleEnd \\
  \text{Natural numbers} &   & \ni & n, m\\
%  \text{Layout contexts} & C & ::= & \many{\ell_i \sim \tau_i} \\
\end{array}
$$
\caption{The syntax of the \Dargent surface language}
\label{fig:dargent-syntax}
\end{figure}


In \autoref{fig:dargent-syntax} we present the surface syntax for \Dargent.
A \emph{layout expression} is a description of the usage of some (heap) memory.
It only describes the low-level view of a memory region---it is not associated
to any particular algebraic data type. From this perspective, \dargent descriptions
are independent of \cogent types. As we will shortly see, however, a given \cogent type can only be laid out in certain ways, which places restrictions on which layouts can be assigned to a given type. 

When specifying a layout, two pieces of information are relevant: How much space a component
occupies in memory and where it is placed in relation to the overall heap object in which it is contained.
Primitive types, such as integer types and
booleans, are laid out as a contiguous block of memory of a particular
size. For example, a contiguous 4-byte block would be an
appropriate layout for the 32-bit word type $\PrimType{U32}$.
A block of memory is specified as a size expression,
which can be in bytes ($\byte$), bits ($\bit$), or additions of smaller sizes.
Additionally, memory blocks of word size (e.g.\ 1 byte, 2 bytes, 4 bytes and 8 bytes)
can be given an endianness ($\bigEnd$ or $\littleEnd$), with the $\using$ keyword.

Components of boxed type are represented as pointers, and thus must be described with the special $\DLPtr$ layout, and not as a chunk of
memory. This special layout improves readability of code, and also allows for some portability:
The $\DLPtr$ layout will have different sizes according to the host machine's architecture.


Layouts for record types use the $\mathsf{record}$ construct, which
contains sub-expressions for the memory layout of each field. As
we can specify memory blocks down to the individual bits, we can naturally
represent records of boolean values as a bit field:
%
\begin{mathpar}
  \mathbf{layout}\ \AbsN{Bitfield} = \DLRec{
     \FieldN{x}: 1\bit
   , \FieldN{y}: 1\bit\ \at\ 1\bit
   , \FieldN{z}: 1\bit\ \at\ 2\bit
  }
\end{mathpar}
%
Here the $\at$ operator is used to place each field at a
different bit offset, so that they do not overlap. If two record fields
have overlapping layouts, the description is rejected by the compiler.

The $\at$ operator can be applied to any layout expressions, which will
shift the entire expression by the specified amount. Alternatively,
the $\after$ operator can be used in a record layout:
\begin{mathpar}
  \mathbf{layout}\ \AbsN{Bitfield} = \DLRec{
     \FieldN{x}: 1\bit
   , \FieldN{y}: 1\bit\ \after\ \FieldN{x}
   , \FieldN{z}: 1\bit\ \after\ \FieldN{y}
  }
\end{mathpar}
so that a later field is placed right after the previous one. This saves
the programmer from calculating the concrete offset. When no offset ($\at$ or $\after$) is given,
it will by default place the field after the previous one.


Layouts for variant types use the $\mathsf{variant}$ construct. It
firstly requires a layout expression for the \emph{tag}. Then, for
each constructor in the variant, a specific tag value needs to be assigned, followed by
a layout expression for the \emph{payload} of that constructor. When one alternative is taken, the memory used by another alternatives becomes irrelevant, which is why the memory for the payloads can overlap.
Additionally, \dargent allows for zero-sized payloads. For instance,
the $\Cons{Maybe}{a}$ type, defined as $\VariantTy{\Cons{Just}{a} \mid \Cons{Nothing}{()}}$,
may be given a layout in which the payload for constructor $\ConsN{Nothing}$ does not occupy
any memory.


% \review{If you encode `Maybe T` as a tag followed by 7 bytes, can you then configure the compiler to replace addition on these `Maybe`s by an optimized computation that leverages the layout choice (e.g. `and`s the flags, performs native byte addition, and then resets the flags)?}
% 
% \ambroise{that could be a nice example of FFI. I'll implement it and describe it if I have time.}

% Array layouts are used for specifying the built-in array types in \Cogent.
% Currently we only allow a uniform layout for all elements of an array, thus
% in the $\mathsf{array}$ layout expression, we only need to supply a layout
% for the element type.

Similar to \cogent types, \dargent expressions are also structural. Layout synonyms
can be defined using the \texttt{layout} keyword, just as
\cogent type synonyms are defined using the \texttt{type} keyword. For example,
%
\begin{mathpar}
\mathbf{layout}\ \AbsN{FourBytes} = 4\byte
\end{mathpar}
%
defines a layout synonym $\AbsN{FourBytes}$, which is definitionally equal to $4\byte$ on
the right hand side. Layout and type synonyms can take parameters.

We can now give a \dargent description to the $\AbsN{Example}$ type in \autoref{fig:cogent-type-ex} (assuming a 64-bit architecture):
$$
\begin{array}{l}
  \mathbf{layout}\ \AbsN{ExampleLayout} = \mathsf{record}\ \{ \\
  \quad \FieldTy{struct}{\mathsf{record}\ \{\colorbox{red!10}{$\FieldTy{x}{4\byte}$}, \colorbox{yellow!30}{$\FieldTy{y}{1\bit}$}\}}, \\
  \quad \colorbox{cyan!10}{$\FieldTy{ptr}{\DLPtr\ \at\ 8\byte}$}, \\
  \quad \FieldTy{sum}{\mathsf{variant}\ (\colorbox{gray!30}{$1\bit$}) \\
  \quad \quad \{\colorbox{blue!10}{$\ConsN{X}(0) : 2\byte\ \at\ 1\byte$}, 
                \colorbox{green!10}{$\ConsN{Y}(1) : 1\byte\ \at\ 1\byte$}}\}\ \at\ 5\byte \\
  \} 
\end{array}
$$

\begin{figure}
\begin{tikzpicture}[background rectangle/.style={fill=white}, show background rectangle,scale=0.80]
  \draw [dashed] (0,0) rectangle (16,2) ;
  \draw [dotted,fill=gray!10] (5,0.1) rectangle (8,1.9) ;
  \draw [fill=red!10] (0.1,0.1) rectangle (4,1.9)node [pos=0.5,align=center] {$\FieldN{x}$};
  \draw [fill=yellow!30] (4.00,0.1) rectangle (4.125,1.9)node [pos=0.5,align=center,anchor=west] {$\FieldN{y}$};
  \draw [fill=gray!30] (5,0.2) rectangle (5.125,1.8);
  \draw [fill=blue!10] (6,1.2) rectangle (8,1.8)node [pos=0.5,align=center] {$\ConsN{X}$};
  \draw [fill=green!10] (6,0.2) rectangle (7,0.8)node [pos=0.5,align=center] {$\ConsN{Y}$};
  \draw [fill=cyan!10] (8,0.1) rectangle (15.9,1.9)node [pos=0.5,align=center] {$\FieldN{ptr}$};
  \draw [->] (5.125,1.5) -- (6,1.5) node [pos=0.5, anchor=south, yshift=-2pt] {\scriptsize if $\!= \!0$};
  \draw [->] (5.125,0.5) -- (6,0.5) node [pos=0.5, anchor=south,yshift=-2pt] {\scriptsize if $= 1$};
  \node [align=center, anchor=north] at (0,0) {0B};
  \node [align=center, anchor=north] at (1,0) {1B};
  \node [align=center, anchor=north] at (2,0) {2B};
  \node [align=center, anchor=north] at (3,0) {3B};
  \node [align=center, anchor=north] at (4,0) {4B};
  \node [align=center, anchor=north] at (5,0) {5B};
  \node [align=center, anchor=north] at (6,0) {6B};
  \node [align=center, anchor=north] at (7,0) {7B};
  \node [align=center, anchor=north] at (8,0) {8B};
  \node [align=center, anchor=north] at (9,0) {9B};
  \node [align=center, anchor=north] at (10,0) {10B};
  \node [align=center, anchor=north] at (11,0) {11B};
  \node [align=center, anchor=north] at (12,0) {12B};
  \node [align=center, anchor=north] at (13,0) {13B};
  \node [align=center, anchor=north] at (14,0) {14B};
  \node [align=center, anchor=north] at (15,0) {15B};
  \node [align=center, anchor=north] at (16,0) {16B};
\end{tikzpicture}
  \caption{The $\mathit{ExampleLayout}$, visualised.}
  \label{fig:example-layout}
\end{figure}
\noindent \autoref{fig:example-layout} gives a pictorial illustration of this memory layout.
In the layout above, it is worth noting that the offsets ($\at\ 1\byte$) for the two payloads of $\ConsN{A}$ and $\ConsN{B}$
are in relation to the beginning of the $\FieldN{sum}$ field, which is $5$ bytes ($5\byte$) from the beginning
of the top-level structure. We can equivalently write $\DLVar{1\bit\ \at\ 5\byte}{\ConsN{A}(0) : 2\byte\ \at\ 6\byte,
\ConsN{B}(1) : 1\byte\ \at\ 6\byte}\ \at\ 0\byte$ for the $\FieldN{sum}$ field without changing its layout.

At this point, this layout is still independent of the \cogent type, and the compiler will only check
that this layout definition is wellformed: that it does not have overlapping fields, that the tag values
are distinct, and so on. To associate a layout to a type, we add a $\layout$ keyword to the type language of \cogent. For this example, the type $\AbsN{Example}$ $\layout$ $\AbsN{ExampleLayout}$ describes the type $\AbsN{Example}$ laid out according to the description in $\AbsN{ExampleLayout}$. The compiler will check that $\AbsN{ExampleLayout}$ is an appropriate layout for the \cogent type $\AbsN{Example}$. We will talk more about the wellformedness and matching rules later
in \autoref{ss:formalised-semantics}. To eliminate verbosity, a type synonym can be given to the layout-annotated type above.


We make the design choice of having \dargent layouts be defined independently 
of \cogent types, and only relating them afterwards with the $\layout$
keyword. This design may seem sub-optimal, as some common information is duplicated
in both the types and the layouts, and matching them requires a set of
dedicated typing rules (as we will
see in \autoref{fig:layout-type}). We make this design decision for several reasons.

Firstly, while the \cogent type structure is central to the functional semantics of any \cogent
program and reasoning about its functional correctness, the exact layout of algebraic data types is just an
implementation detail, totally irrelevant to the top-level Isabelle/HOL
embedding of the \cogent program, and whose correctness is guaranteed by the automatic
C refinement proof that the \cogent compiler produces. 
%
Thus, type structure and layout are conceptually separate.
%These two constructs are conceptually separate.

Secondly, this approach is more flexible and amenable to future
extensions. Currently, as we will see in \autoref{fig:layout-type}, the
layout--type matching is fairly restricted: e.g.\ a $\PrimType{U16}$
type has to occupy 2 bytes ($2\byte$) in memory, and a record type has
to be laid out in accordance with a $\mathsf{record}$ layout.  In the
future, several extensions could be made to relax this matching
relation. For example, it is technically valid to store a smaller type
in a larger memory area, say, a $\PrimType{U16}$ type in a memory region
of 4 bytes, or a record type in a contiguous chunk of memory that is big
enough. Some heuristics can be implemented in the compiler to decide how
to arrange underspecified layouts. This feature can be useful in
improving programmers' productivity, permitting the omission of layout
details for unimportant parts of a type. Also, as there is ongoing work
towards adding refinement types to \cogent~\citep{Paradeza_20}, a
refined type could possibly be laid out in a smaller memory area. For
instance, $\{ \nu : \PrimType{U16}\ |\ \nu < 2 \}$ can be laid out in a
memory area as small as 1 bit. None of these extensions would be easy to
implement if the layout information was baked into the types.

Thirdly, separating layouts and types encourages modularity.
Developers using the \cogent language can write programs without
needing to worry about the detailed layout of types,
and are still able to prove functional correctness of their code
against some high-level specification and ship their code to end users.
The end users, with the knowledge of the particular target architecture and environment,
can decide the layouts and plug them into the \cogent programs.

Finally, although our design requires a dedicated set of rules for
checking the layout--type matching, it significantly simplifies compiler
engineering in the long term. The \cogent compiler is very large, and
the layout--type matching checker only constitutes a small part of the
typechecker. The compiler not only compiles \cogent programs to C, but
also generates information for various Isabelle proof tactics, numerous
embeddings of the program in Isabelle and Haskell~\citep{Chen:22}. A lot
of these embeddings are only concerned with types, and not layouts. If
the layouts and types were merged, any changes to the layout
implementation would require changes to various irrelevant parts of the
compiler.


\subsection{Layout Polymorphism}

% \dargent layouts are only concerned of low-level implementation details of a \cogent program,
% and they do not affect the functional semantics of a program. It means that program semantics
% is parametric over layouts.
We extend \Cogent's existing parametric polymorphism mechanism to support \emph{layout polymorphism}. This feature allows users to abstract
over the layout that they would like to assign to a certain type.
In a \cogent function signature, \emph{layout variables}, just as type variables, can
be universally quantified. These layout variables may be constrained, similarly to type constraints in Haskell, which require that a layout variable matches a type.\footnote{A formal definition of layout matching is discussed in \autoref{ss:formalised-semantics}.}
For example, in the code snippet below,
\begin{align*}
  & \mathbf{type}\ \Cons{Pair}{t} = \RecordTy{\FieldTy{fst}{\VarN{t}}, \FieldTy{snd}{\VarN{t}}} \\
  & \layout\ \Cons{LPair}{l} = \DLRec{\FieldTy{fst}{\VarN{l}}, \FieldTy{snd}{\VarN{l}\ \at\ 4\byte}}\\
  & \FunN{freePair} : \forall (t, l :\sim t).\; \Cons{Pair}{t}\ \withL{\Cons{LPair}{l}} \rightarrow ()
\end{align*}
we define a parametric type synonym $\Cons{Pair}{t}$ and layout synonym $\Cons{LPair}{l}$, and
an abstract polymorphic function that operates on such a pair. In the function's type, we require layout $l$ to be compatible with type $t$, so that the $\Cons{LPair}{l}$
is always a valid layout expression for type $\Cons{Pair}{t}$.
Layout-polymorphic functions may be explicitly applied to layouts, akin to explicit type applications.
For example, $\TyLyApp{freePair}{\PrimType{U32}}{\AbsN{FourBytes}}$ instantiates $t$ to $\PrimType{U32}$
and $l$ to $\AbsN{FourBytes}$. If the type/layout application is
incompatible, as for instance in $\TyLyApp{freePair}{\PrimType{U8}}{1\bit}$, where
$1\bit$ is not big enough to store the $\PrimType{U8}$, the typechecker will reject
such a program.  The typechecker also ensures that any instantiation of $l$ produces layouts
that are wellformed. For example, if $l$ is instantiated with $8\byte$, it will
render $\Cons{LPair}{l}$ ill-formed, as the $\FieldN{snd}$ field will overlap with the
$\FieldN{fst}$. This can be rectified by using the $\after$ relative location (or
leaving the location implicit)
instead, which will automatically place the $\FieldN{snd}$ field right after $\FieldN{fst}$.

Layout polymorphism is necessary to retain the full generality of type
polymorphism in the presence of \dargent, as demonstrated by the example
above.
%
Layout polymorphism also facilitates code reuse in several scenarios:
\begin{enumerate*}
  \item It can be used in programs that are architecture-agnostic.
    For the same program running on different architectures, the same algebraic datatype
    may be laid out differently according to the hardware it runs on. These datatypes are
    not necessarily part of the API, as developers tend to design interfaces
    in an architecture-independent manner and these types thus have uniform layouts
    across different architectures. However, datatypes that are internal to the program
    can be represented with appropriate layouts to best suit the architecture and the
    respective C compiler for better performance.
  \item Besides architectural differences, layout polymorphism allows types to be reused more
    broadly across different applications.
    For instance, colours are a general concept independent of the application. One common
    way to logically represent a colour is the ARGB model. However, the actual low-level layouts
    for the logical ARGB value vary. Commonly used layouts include ARGB32 and RGBA32, and
    other layouts also exist. In this case, developers can simply define an ARGB colour type with a polymorphic layout to be determined at use-sites.
  \item Layout polymorphism also facilitates software engineering practices.
    For example, developers can start with the default layout without a \dargent annotation,
    and \emph{incrementally}
    optimise the program to use more compact and efficient \dargent layouts. In this
    case, they simply instantiate the layouts differently, without needing to
    re-define the types. Benchmarking is another use case: developers can easily have
    the same type with different instantiations \emph{side-by-side} to study their performance
    differences.
\end{enumerate*}


%In Section~\ref{ssec:bitfields}, we present another extension that accounts for custom sized integers such as $\ConsN{U7}$ or $\ConsN{U15}$. \liam{Don't understand why this is mentioned here}

%Additionally, choosing to keep layouts and types separate does not preclude us from providing users with syntactic sugar that eliminates information duplication. 
% implementing syntactic sugar
%to eliminate information duplication
% and improve the user
%experience of the language. 
%We leave it
%as future work.

\subsection{The \Dargent Core Calculus}\label{ss:dargent-core}

\begin{figure}
\[
\begin{array}{llclr}
  \text{Bit ranges} & r\\
  \text{Offsets}    & o & \in & \mathbb{N} \\
  \text{Layouts} & \ell & ::=  & () & \text{(unit layout)}\\ 
                 &      & \mid & r_\omega & \text{(primitive layout)} \\
                 &      & \mid & x\{o\} & \text{(layout variable)} \\
                 &      & \mid & \multicolumn{2}{l}{\mathsf{record}\ \{ \overline{\mathrm{f}_i\!: \ell_i} \}}\\
                 &      & \mid & \multicolumn{2}{l}{\mathsf{variant}\ (\ell)\ \{ \overline{\mathrm{A}_i\ (n_i)\!: \ell_i} \}}\\
 %                &      & \mid & \multicolumn{2}{l}{\mathsf{array}\ \{ \ell \}}\\
  \text{Field names} & & \ni & \mathrm{f}\\
  \text{Constructors} & & \ni & \mathrm{A}\\
  \text{Endianness} & \omega & ::= & \multicolumn{2}{l}{\mathsf{BE} \mid \mathsf{LE} \mid \mathsf{ME}} \\
%  \text{Layout contexts} & C & ::= & \many{\ell_i \sim \tau_i} \\
\end{array}
\]
\caption{The syntax of the \Dargent core language}
\label{fig:dargent-core}
\end{figure}




The \Dargent surface language is desugared into a smaller core calculus,
whose syntax is outlined in \autoref{fig:dargent-core}. The core calculus is
the language on which the verification is based. As can be seen in  \autoref{fig:dargent-core}, layouts consist fundamentally of \emph{bit ranges}, which describe which bits in memory are used
to store each piece of data. The definition of bit ranges is omitted, because our core calculus is parametric over the exact bit range
representation---at different points in the compilation pipeline, bit ranges are represented
differently, to suit the purpose of that particular phase of the compilation.

Bit ranges are annotated with an endianness to form a primitive layout in our 
core language. When the endianness is not specified in the surface language, 
a \emph{machine endianness} will be given by default, written $\mathsf{ME}$
 in core. When C code is generated, a
subroutine will be invoked to determine the machine endianness, so that the C code works
as intended.
%
In this paper, when the endianness is unimportant, we will omit the
subscript from $r_{\omega}$.

When the surface layout expressions are first desugared, the bit ranges are represented
as $\bitrange{o}{s}$, which is a pair of natural numbers, indicating the offset ($o$ bits) from
\emph{the beginning of the top-level heap object} in which it is contained, 
as well as the range occupied ($s$ bits).
%
In $\bitrange{o}{s}$, we require $s > 0$; zero-sized layouts can use the
empty layout $()$ instead.

Layout variables are given the form $x\{o\}$, since we must remember to
apply the offset after instantiation---when $\VarN{x}$
is instantiated, it will be shifted to the right by $\VarN{o}$ bits.
The other core layouts are very similar to their counterparts in the
surface language.

In summary, the desugarer converting the surface layout expressions to
the core calculus must perform the following tasks:
\begin{itemize}
  \item expand layout names to layout definitions;
  \item translate size expressions into bit ranges;
  \item compute offsets relative to the beginning of the top-level
    heap object;
  \item insert explicit layout applications to any layout-polymorphic
    function calls.
\end{itemize}
% The desugaring algorithm is a simple inductive definition, which can be 
% found in~\autoref{ss:static-semantics}.
\Dargent is used for describing the layout of heap memory. To cover all the
heap memory that is addressable from an object,
it suffices to attach a layout description to each pointer that the object contains.
In the $\AbsN{Example}$ type that we showed earlier in \autoref{ssec:informal-intro},
the layout annotation $\AbsN{ExampleLayout}$ is attached to the $\AbsN{Example}$ type,
which is a boxed type. That layout description contains all the information about
how the fields should be laid out, in particular the unboxed fields $\FieldN{struct}$ and
$\FieldN{sum}$ which also reside in the heap. For the boxed field $\FieldN{ptr}$, however,
the description in $\AbsN{ExampleLayout}$ does not extend beyond the indirection---it only
knows that $\FieldN{ptr}$ is a $\DLPtr$ but it is oblivious to the
memory layout behind $\FieldN{ptr}$.
%
To prescribe the layout for $\{ \FieldN{c} : \PrimType{U8}\}$, a
separate layout annotation should be attached to this type, \emph{e.g.}
$\{ \FieldTy{c}{\PrimType{U8}} \}\ \layout\ \DLRec{
  \FieldTy{c}{1\byte}}$.

Recall that, if an object is referenced by-pointer, it has a boxed type.
This means that in the core calculus, data layouts only need to be attached to boxed sigils
and nothing else.
We modify the definition of sigils as follows:
$$
\begin{array}{llclr}
  \text{Sigils} & s & ::=  & \Writable{\ell} \mid \ReadOnly{\ell} \mid \Unboxed &
\end{array}
$$
If the sigil is boxed (either writable or read-only), it can carry a \dargent layout.
We use $\Boxed{\ell}$ as a notational convenience when we do not wish
to distinguish its mutability.
When the layout annotations are left out, the type will be compiled with the default layout chosen
by the compiler.



\subsection{The Static Semantics with \Dargent} \label{ss:formalised-semantics}

% \review{any proof of the soundness of the type system?}

%type system, kinding system for Cogent
%changes needed for Dargent both to type and kinding system
% we need to ensure that types and layouts match.
% Since cogent has polymorphic types and layouts. 
% We do not always have sufficient information at the type checking phase to check that types and layouts match. 
% we therefore need to track a set of pairs of types and layouts that we need to check matching for at a later stage of compilation (after monomorphisation).

% lemma 3 in cogent paper 
%Lemma 3 (Type instantiation preserves kinds). For any type τ, αi:Kκi ⊢τ:K κimplies∆⊢τ[ρi/αi]:K κwhen,foreachi, ∆⊢ρi :K κi.
% would no longer hold once we introduce layouts
%why? because now typing and kinding also need to substitute layouts and 
% since typing should match layouts and types, we need to track type variables and their layouts to ensure once they are instantiated, their instatiation matches the given layouts. Layouts can also be polymorphic and the type-layout matching should also be respected when layouts are instantiated. 
%
%Since layouts are also polymorphic, we also need to instantiate them
%1) if substitution is wf and type is wf then resulting type is wf. (why? subject reduction?)
%% need to introduce a set of constraints C

% 2) if a type is a subtype (in a linear sense) of a wf type then it is also wf and vice-versa.
% now a type wfness includes constraints C
% any type that includes a polym. type variable or poly layout needs to be checked for wf later on once these are instantiated. We therefore need to keep track of such types and layouts in order tpo check later on that they match and respect layout wf.  
% suppose you have a record with a layout, the context must contain record type with layout for the layout to be wf. 
% record type t with layout l (C must contain pair (t, l) )
% now we want to know that (bang(t), l) is also wellformed (why?)
% so constraints C must contain (bang (t), l)
%  linearity of types does not affect their layout in memory.
% i.e. whether types are read-only or writable
%therefore the matching between types and layouts can use simplified types (with linearity info removed)
% in fact this is necessary, why?
%  any type (regardless of linearity ) relates to layout
% because we want 3
% 3) if a type t is wf then the bang (t) is wellformed (turning to non-linear read only version)
% \delta = [\tau_1.. tau_n]
% [\tau:= \delta \tau], 

%Thereofore, Dargent's typing judgement needs to therefore additionally track 

% can't do that directly but rather need to keep track of a set of pairs (l,t) where t is a simplified type.
%
%a set of (l,t), the typesystem needs to ensure they match
% why simplified ?


%This section, we give a formal account of the static semantics of the core
%calculus on which we base our \dargent formalisation.

\Cogent's static semantics, along with some meta-properties
of the type system such as type preservation, are formalised in Isabelle/HOL
(\citet{OConnor:2016,Cogent_JFP_2021} give more details). In this section, we explain
how these formal proofs and definitions are extended to account for \Dargent.

In the original \Cogent typing judgement $\Delta;\Gamma \vdash e : \tau$, we have the set of polymorphic type variables $\Delta$ in the context. Each entry in $\Delta$ also includes a set of \emph{constraints} for that type variable, which may enforce, for example, a type variable represent only types that are shareable.
 When these type variables are instantiated, these constraints are checked. We also have a wellformedness judgement on types, $\Delta \vdash \tau\ \mathbf{wf}$, which ensures that all type variables are in scope in $\Delta$ and that the type structure of $\tau$ respects the linearity constraints in $\Delta$. 
%Given a \cogent type, \TODO{say something about subtyping}.  
%\Cogent's kinding judgement, $\Delta \vdash \tau :_K \kappa$, ensures that \cogent types are wellformed with respect to linearity constraints specified in the sigil.  
%$\Delta \vdash \tau :_K \kappa$, stating that the type $\tau$ is wellformed
%in the context of type variables $\Delta$. 

With \Dargent, the typing judgement must additionally ensure that \cogent types
match their assigned layouts (\autoref{fig:layout-type}) and that the layouts
are wellformed (\autoref{fig:wf-layout}). Because \cogent has
polymorphic type and layout variables, however, these additional constraints
cannot always be immediately ensured when typechecking locally. Thus, we
extend the \cogent typing judgement $\Delta;C;\Gamma \vdash e : \tau$ and
wellformedness judgement $\Delta; C \vdash \tau\ \mathbf{wf}$ to include an
additional set of \emph{layout} constraints $C$. It
is a set of pairs $(\ell,\hat{\tau})$ denoting that the layout $\ell$ must
be capable of representing type $\hat{\tau}$. When the types and layouts are instantiated, these matching constraints are checked.


\begin{figure}
\[
\begin{array}{llclr}
  \text{Simplified types} 
     & \hat{\tau} & ::=  & T \mid \mathsf{pointer} \mid () \mid \alpha \mid
                         \{ \many{\mathrm{f}_i : \hat{\tau}_i} \} \mid
                         \VariantTy{\many{\mathrm{A}_i\ \hat{\tau}_i}} \mid \dots &
\end{array}
\]
\caption{The syntax of simplified types}
\label{fig:syntax-simplified-types}
\end{figure}


The type in the context $C$ is \emph{simplified}, denoted as $\hat{\tau}$,
because any type structure beyond the single heap object described by the
layout $\ell$ is irrelevant. This means that all boxed types, which are all
uniformly represented as a pointer to another heap object, can be treated
simply as a primitive machine word when matching with layouts. The syntax of (a
subset of) simplified types is given in \autoref{fig:syntax-simplified-types}
(cf.\ \autoref{fig:syntax-cogent-types}). Types are simplified by
replacing all boxed components of the type with the newly introduced
$\mathsf{pointer}$ type constant.  

%
%
%
%Two important things to note at this stage: 
%(1) \cogent types can contain polymorphic type variables and \dargent supports polymorphic layouts. 
%(2) the layout of a type is independent of its linearity constraints
%
%In particular, the typing judgement in the presence of \dargent $\Delta;C;\Gamma\vdash e : \tau$ additionally track a set of pairs $C$ of layouts $\ell$ and \emph{simplified types} $\hat{t}$ (\autoref{fig:syntax-simplified-types}). \dargent's type system  maintains matching constraints between the types and the layouts. A simplified type is just a \Cogent type with the linearity information removed. 

%In this section, we give a formal account of the static semantics of the core
%calculus on which we base our \dargent formalisation.
%The typing rules of \dargent differ from those of the \Cogent  in that they involve an additional environment that lists constraints on relating a layout to a \emph{simplified type}
%(\autoref{fig:syntax-simplified-types}).


%A simplified type %
%\footnote{The notion of simplified types introduced here is very similar
%  to the \emph{representation types} used in the C refinement proof. The main difference is that
%  it now includes type variables, whereas the representation types don't, which makes sense as
%  C is not polymorphic. The properties that can be proved on this simplified types are also very similar
%  to those of representation types, especially stabilty under subtyping, which is a main
%  motivation and contribution of the subtyping design.}
%is just a \cogent type with all the linearity information removed. As such, the simplified type is stable under
%subtyping and $\mathbf{bang}(\cdot)$. 
%(turning a linear type into a read-only
%non-linear one), thus making constraints invariant
%under these operations. 
%These properties are used when adapting the proofs that subtyping and $\mathbf{bang}(\cdot)$ preserve wellformedness of types. 
%More precisely, the function
%$\textbf{simplify}$ that turns a regular \cogent type into a simplified type (boxed
%types are turned into pointers). Then,
%$\textbf{simplify}(\BangF{\tau}) = \textbf{simplify}(\tau) $ and
%$\textbf{simplify}(\tau_1) = \textbf{simplify}(\tau_2)$ if $\tau_1$ is a subtype
%of $\tau_2$.

\subsubsection{Typing Rule for Records with Layouts}

\newcommand{\highlight}[1]{\colorbox{blue!10}{$\displaystyle #1$}}
\newcommand{\highlighttext}[1]{\colorbox{blue!10}{#1}}

We are now ready to present the main typing rule that is added for
\dargent: wellformedness of boxed records with custom layouts.
\begin{equation*}
  \label{eq:wellformed-rec}
   \inferrule
	{\overline{\mathrm{f}_i}\ \text{distinct} \\
   \text{for each $i$: }\Delta; C \vdash \tau_i\ \mathbf{wf} \\\\
     \highlight{\ell\ \mathbf{wf}} \\
     \highlight{
       % \text{if $\ell$ or some $\tau_i$ is open, then}
       \ell\sim \{ \overline{\mathrm{f}_i :
         \hat\tau_i}\} }  
     \\
	 \highlight{
       % \text{if $\ell$ or some $\tau_i$ is open, then}
       (\ell,\{ \overline{\mathrm{f}_i :
         \hat\tau_i}\}) \in C}
   }
	{\Delta ; C\vdash \{ \overline{\mathrm{f}_i :: \tau_i} \} b_\ell\ \mathbf{wf}}
   {}
  {\textsc{RecWf}}
\end{equation*}
The first two premises in this rule are similar to those in \cogent, the \highlighttext{next three premises} are added to account for \dargent. The layout wellformedness judgement, $\ell\ \mathbf{wf}$, ensures that field layouts do not overlap (\autoref{fig:wf-layout}). 
The last two premises require the type to match the layout (\autoref{fig:layout-type}) and ensure that this is tracked in the set of constraints, respectively.
% \footnote{
%   The notion of context does not enforce that if $(\ell, \hat{t}) \in C$,
%   then $\ell \sim \hat{t}$, but 
%   \textsc{RecWf} is the only rule that requires some
%   constraint to be in $C$, and it indeed enforces it.
% }% 
In practice, we use a variant of the rule \textsc{RecWf} where the
layout constraint is only tracked in $C$ if the type or the layout are
polymorphic.



%These two premises are necessary for ensuring that any valid instantiation of a well-typed polymorphic program is well-typed which in turn is needed for proving type preservation. 
%the the following specialisation lemmas hold.
\begin{figure}
\begin{inductive}{\ell \sim \hat{\tau}}
  \inferrule
    {\mathrm{bits}\;(T) = s}
    {\bitrange{o}{s} \sim T}
    {\textsc{PrimTyMatch}}
  \qquad
  \begin{array}{ll}
    \bits{\ConsN{U}n}   & = n  \\
    \bits{\ConsN{Bool}} & = 1
  \end{array}
  \\
  \inferrule
    {M\ \text{is the pointer size}}
    {\bitrange{o}{M} \sim \mathsf{pointer}}
    {\textsc{BoxedTyMatch}}
  \qquad
  \inferrule
    { }
    {x\{o\} \sim \hat{\tau}}
    {\textsc{VarMatch}}
  \\
  \inferrule
    { }
    {() \sim ()}
    {\textsc{UnitMatch}}
  \qquad
  \inferrule
    {\text{for each $i$: } \ell_i \sim \hat{\tau}_i}
    {\DLRec{\many{\FieldN{f}_i : \ell_i}} \sim \{\many{\FieldN{f}_i : \hat{\tau}_i}\}}
    {\textsc{URecordMatch}}
  % \\
  % \inferrule
  %   {\ell \sim \tau \\
  %    \ell\ \text{is byte aligned}}
  %   {\DLArr{\ell} \sim \ArrayTy{\tau}{v}{\Unboxed}}
  %   {\textsc{UArrayMatch}}
  \\
  \inferrule
    {\text{for each $i$:\ } \ell_i \sim \hat{\tau}_i}
    {\DLVar{s}{\many{\ConsN{A}_i\ (v_i) : \ell_i}} \sim \VariantTy{\many{\ConsN{A}_i\ \hat{\tau}_i}}}
    {\textsc{VariantMatch}}
\end{inductive}
\caption{Matching relation between layouts and types}
\label{fig:layout-type}
\end{figure}



\begin{figure}
\begin{inductive}{\ell\ \mathbf{wf}}
  \inferrule
    { }
    {\bitrange{o}{s}\ \mathbf{wf}}
    {\textsc{BitRangeWF}}
    \qquad
    \inferrule
    { }
    {()\ \mathbf{wf}}
    {\textsc{UnitWF}}
    \qquad
    \inferrule
    { }
    {x\{o\}\ \mathbf{wf}}
    {\textsc{VarWF}}
  \\
  \inferrule
  {\text{for each $i$: } \ell_i \ \mathbf{wf} \\\\
    \text{for each $i\neq j$: }
    \mathrm{taken}(\ell_i) \cap 
    \mathrm{taken}(\ell_j) = \emptyset
  }
    {\DLRec{\many{\FieldN{f}_i : \ell_i}} \  \mathbf{wf}}
    {\textsc{URecordWF}}
  % \\
  % \inferrule
  %   {\ell \sim \tau \\
  %    \ell\ \text{is byte aligned}}
  %   {\DLArr{\ell} \sim \ArrayTy{\tau}{v}{\Unboxed}}
  %   {\textsc{UArrayWF}}
  \\
  \inferrule
  {\text{for each $i$:\ } \ell_i \ \mathbf{wf}\\
    v_i < 2^s
    \\
    \mathrm{taken}(\ell_i) \cap \mathrm{taken}(\bitrange{o}{s}) = \emptyset
    \\\\
    \text{for each $i\neq j$:\ } v_i \neq v_j }
    {\DLVar{\bitrange{o}{s}}{\many{\ConsN{A}_i\ (v_i) : \ell_i}}\ \mathbf{wf}}
    {\textsc{VariantWF}}
\end{inductive}

  where
  $\mathrm{taken}(\ell) \in \mathbb{N}$
   returns the set of bit positions that $\ell$ occupies
    (defined recursively).
\caption{Wellformed layouts}
\label{fig:wf-layout}
\end{figure}




\subsubsection{Specialisation}

% change substitution syntax to something more readable, remove kinding and just keep wellformedness of types if possible
\cogent's type preservation proof relies on type specialisation maintaining welltypedness~\citep[Lemmas 3 and 4]{OConnor:2016}. 
With \Dargent, we need to adapt these statements as follows.
%We state them below and discuss them afterwards, forgetting kinding information in $\Delta,\Delta'$ for the sake of simplicity.
\begin{lemma}[Specialisation lemmas]
  \label{lem:specialisation-lemmas}
  Let $\overline{\rho_i}$ be a list of types and $\overline{\ell_j}$ be a list
  of layouts. Let $\overline{\alpha_i}$ denote the list of type variables declared
  in $\Delta$.

  Then, $\Delta;C \vdash \tau\ \mathbf{wf} $ and  $\Delta;C;\Gamma \vdash e : \tau$  
  respectively imply 
     \begin{align*}
      \Delta';C'& \vdash \tau[\overline{\rho_i}/\overline{\alpha_i}, \overline{\ell_j}/\overline{x_j}]\ \mathbf{wf}; \\
      \Delta';C';\Gamma[\overline{\rho_i}/\overline{\alpha_i}, \overline{\ell_j}/\overline{x_j}]%^\dagger 
      & \vdash e[\overline{\rho_i}/\overline{\alpha_i}, \overline{\ell_j}/\overline{x_j}]: \tau[\overline{\rho_i}/\overline{\alpha_i}, \overline{\ell_j}/\overline{x_j}  ],
     \end{align*}
  under the following conditions:
  \begin{itemize}
    \item for each $i$, $\Delta', C'\vdash \rho_i\ \textbf{wf}$;
    \item for each pair $(\ell, \hat{\tau}) \in C$ such that
      $\Delta \vdash \hat{\tau}\ \mathbf{wf}$, the following statements hold:\footnote{The
        condition $\Delta \vdash \hat{\tau}\ \mathbf{wf}$ is required to support a stronger induction hypothesis. As always, the gory details are in the machine-checked proofs. }
      \begin{itemize}
      \item $\ell[\overline{\ell_j}/\overline{x_j}]\ \mathbf{wf}$;
      \item 
        $ \ell[\overline{\ell_j}/\overline{x_j}] \sim \hat{\tau}[\overline{\hat{\rho}_i}/\overline{\alpha_i}]$;
       \item $(\ell[\overline{\ell_j}/\overline{x_j}], \hat{\tau}[\overline{\hat{\rho}_i}/\overline{\alpha_i}]) \in C' $.
      \end{itemize}
  \end{itemize}
\end{lemma}

% \begin{lemma}[Type wellformedness specialisation]
%   \label{lem:type-specialisation}
%   If $\Delta;C \vdash \tau\ \mathbf{wf} $ and
%   $\Delta';C' \vdash \overline{\rho_i},\overline{l_j} :_s \Delta; C$, then
%   $$\Delta';C' \vdash \tau[\overline{\rho_i}/\overline{\alpha_i}, \overline{l_j}/\overline{x_j}] \mathbf{wf} $$
% \end{lemma}

\noindent In this lemma,  $[\overline{\rho_i}/\overline{\alpha_i}, \overline{\ell_j}/\overline{x_j}]$ denotes the simultaneous
substitution replacing both type variables $\alpha_i$ and layout variables $x_j$,
and $\Gamma$ is a typing context environment.


Informally, the two conditions mean that the pair 
$(\overline{\rho_i}, \overline{\ell_j})$ of type and layout lists defines a valid 
substitution from $\Delta;C$ to $\Delta';C'$. %, taking $\Delta := \overline{\alpha_i}$,
They also appear in the typing rule for specialising polymorphic functions.
The second one is new and enforces that for each  pair
$(\ell, \hat{\tau})\in C$ that is wellformed (in the sense that type variables
appearing in $\hat{\tau}$ are in $\Delta$), 
the substituted constraint is satisfied and remembered in the set of constraints $C'$.
This last requirement can in fact be dropped if the substituted constraint is closed.
%   Formally, the substitution wellformedness rule is defined as:

% $$
% \inferrule
%   { \text{for each}\ \kappa_i\ \text{in}\ \Delta\ \text{and}\ t_i\ \text{in}\ \delta:\ \Delta'; C' \vdash t_i :_K \kappa_i \\\\
%     \text{for each}\ (\ell_j, \hat{t}_j)\ \text{in}\ C:\ \ell_j[\varsigma] \sim \hat{t_j}[\textbf{simplify} \circ \delta] \\\\
%     (\ell_j[\varsigma], \hat{t_j}[\textbf{simplify}\circ \delta])\in C'
%   }
%   {\Delta'; C' \vdash \varsigma, \delta :_s \Delta; C }
%   {\textsc{SubstWf}}
% $$


% A similar specialisation lemma holds for well-typed terms as well:
% \begin{lemma}[Typing specialisation]
%   \label{lem:typing-specialisation}
%   If $\Delta;C;\Gamma \vdash e : t$ and
%   $\Delta';C' \vdash \varsigma,\delta :_s \Delta; C$, then
% 	$$\Delta';C';\Gamma[\delta, \varsigma]^\dagger \vdash e[\delta, \varsigma]: t[\delta, \varsigma].
% 	~\footnote{$\cdot[\cdot,\cdot]^\dagger$ is the simultanous substitution lifted to the environment.}
%   $$
% \end{lemma}

% Now, let us explain why the wellformedness rule \textsc{RecWf} for records
% not only requires a check that the layout matches the type, but also that
% this constraint is remembered. Indeed, because the type and the layout may
% contain variables, this matching is only partial and would relate
% the layout $ \mathsf{record}\ \{ a : x_1, b : x_2\}$ to
% the type $ \mathsf{record}\ \{ a : y_1, b : y_2\} $.
% But a substitution may later specialise $x_{1/2}$ and $y_{1/2}$ to incompatible
% layouts and types in various ways (including the overlapping of layouts), hence the need for remembering the constraint and checking it again
% when encountering a substitution.


% This rule is actually the only one that requires a constraint to be present in $C$.
% Note that such a constraint always relates a layout with a record type.

% The typing rule for specialising polymorphic functions now uses the
% specific notion of wellformedness for substitutions:
% \[
%   \inferrule
%   {
%     \FunTyEnv{f}{\forall \Delta'.C'
%       \Rightarrow \tau \rightarrow \tau'}  \\
%     {\Delta;C \vdash \varsigma,\delta :_s \Delta'; C'}
%   }
%   {\Delta ; C ;\Gamma \vdash f[\delta,\varsigma] : (\tau \rightarrow \tau')[\delta,\varsigma] }
%   {\textsc{Fun}}
% \]


%\TODO{The subtyping needs to be adapted to allow reordering the fields
%(or not, as we may assume they are already canonically ordered, since the
%formalisation supposedly accounts for the core syntax?)}

%\subsubsection{Future extensions}
%% Introducing this general set of constraints in the context could be avoided
%% if we restrict specialising layout variables to monomorphised types.
%
%% In this case, we could also reuse the already formalised simplified types
%% which do not have any variable case.
%
%% By exploiting the preservation of the typing rules by increasing the set of
%% constraints, it is straightforward to design a procedure that infers the necessary
%% set of constraints for a layout polymorphic function to be well-typed (if it can
%% be).
%
%Our type system currently does not check that the constraints
%in $C$ are compatible. For example, nothing forbids $C$ from containing both
%$(x\{o\},\ConsN{U8})$ and $(x\{o\}, \ConsN{U16})$, although in that case the layout variable 
%$x$ can never be fully instantiated, because of the wellformed rule for 
%substitution mentioned above.
%%
%Note that this does not allow unsafe behaviour, since C code generation only
%involves completely monomorphised layouts and types.
%%
%We leave for future work the design of a (possibly partial) validator that would
%deal with such situations and could raise a typing error as early as possible before
%instantiation.


\subsection{Compiling Records: Custom Getters and Setters}
\label{ssec:gsetters}

Currently, records are the only built-in types
in \cogent that can be boxed. We therefore use records as an example to
discuss how \dargent layouts influence the generated C code.

Without custom layouts, a \cogent record is directly compiled to a C struct
with as many fields as the original record: If
$T = \{a : A, b : B\}$, then
$\sem{ T }$ is a pointer type to
$\texttt{struct}\ \{ \sem{ A }\ a; \sem{ B }\ b; \}$,
 where $\sem{ \cdot }$ denotes the compilation of \cogent types to C types.


The \Dargent extension to the compiler relies on the observation that
we are free to choose
what a \cogent boxed record is compiled to
as long as we provide getters and setters for each field, since they
account for all the available \cogent operations on boxed records.
Assigning a custom layout
$\ell$ to a \cogent record $T$ results in a C type that we denote by
$\sem{ T }_\ell$, consisting of a C struct with a fixed-sized array
of words as a single field.
The implementation chooses the word size to be 32 bits, primarily because
most \cogent programs we develop are targeting 32-bit embedded systems, but this is not fundamental to the design and can be easily made 
configurable to any word size.
It is worth mentioning that this does not mean that a layout has to be a neat
multiple of 32 bits in size. It is absolutely valid to have a record
layout like $\DLRec{ \FieldTy{a}{1\byte}, \FieldTy{b}{3\bit} }$, 
11 bits in total. When this layout is embedded in another layout,
e.g.\ $\DLRec{ \FieldTy{x}{\DLRec{ \FieldTy{a}{1\byte}, \FieldTy{b}{3\bit} }},
               \FieldTy{y}{2\byte} }$, 
the remaining bits after the field $\FieldN{b}$ will be used by
the following field $\FieldN{y}$, without any implicit padding.


The getters and setters for each field are generated
according to the layout. If a top-level boxed record $T$ contains
a field $a : A$, the C prototypes are
\begin{align*}
    & \sem{ A } \texttt{ get\_a }(\sem{ T }_\ell \; t) ;
\\
    & \texttt{void set\_a }(\sem{ T }_\ell  \; t, \sem{ A } \; a) ;
\end{align*}
Note that $\sem{A}$, the return type of the getter (and similarly for
the second argument of the setter function) does not involve the
layout $\ell$. If $A$ is a boxed type, then $\sem{A}$ is a pointer to the
type $A$, whose layout is dictated by the layout information stored in
$A$'s sigil, independent of $\ell$. If $A$ is unboxed, the getter function is the point at which
we convert the low-level custom representation governed by the layout $\ell$
into a standard representation of $A$, so that the value of the field
$a$ can be inspected by the rest of the program. 
As an example, consider the $\FieldN{struct}$ field in \autoref{fig:cogent-type-ex}.
Roughly, the generated C getter has prototype
\cc|struct { U32 a; bool b; } get_struct (Example * t)|, where
$\AbsN{Example}$ is a C structure with a single field consisting of a fixed-sized array
spanning 16~bytes.
% Note that accessing the nested subfield
% \cg|a| or \cg{b} in \cogent involves first retrieving the whole top-level field
% \cg|struct|, which is inefficient and could be improved in future
% work, in particular since nested field getters and setters are already generated
% anyway, as auxiliary functions.


% \TODO{zilinc: describe the nested getters/setters more carefully
% to avoid the confusion that I had.}

% \subsubsection{Implementation}
% \review{Meanwhile, the paper spends a lot of space talking about things like layout
%   variables that seem neither very surprising nor very important. That space would
%   be used better to discuss, for example, the details of the generated getters and
%   setters.}

Getters and setters are generated recursively following the structure of the
involved field types. The process typically involves generating
auxiliary ``nested setters and getters'' that
directly manipulate the data array based on the value of the nested field (such as
$\FieldN{a}$ or $\FieldN{b}$ in the example of \autoref{fig:cogent-type-ex}), and similarly
for getters.
% For example, For nested unboxed records, such as \cg|struct| in
% Figure~\ref{fig:cogent-type-ex}, the compiler generates auxiliary ``nested
% setters and getters'' that
% directly manipulate the array based on the value of the nested field (such as
% $\FieldN{a}$ or $\FieldN{b}$), and similarly
% for getters.
 % Setting a field which is an unboxed record results in successive calls to these
 % nested setters.
For example, the getter and setter for the \cg|struct| field of 
\autoref{fig:cogent-type-ex} are roughly implemented as follows:
\begin{lstlisting}[language=C]
  // the data array
  typedef struct Example { U32[4] data } Example;

  struct field_struct { U32 a; bool b; };

  // setter for the struct field
  void set_struct(Example * d, struct field_struct v){
    // calls to nested setters
    set_struct_a(d, v.a);
    set_struct_b(d, v.b);
  } 

  // getter for the struct field
  struct field_struct get_struct(Example * d){
    // calls to nested getters
    return (struct field_struct){.a = get_struct_a(d), .b = get_struct_b(d)};
  } 
 \end{lstlisting}
% Getters follow a similar pattern.
% As for variant field, such as 
% Similarly, updating a variant field results in first updating the tag value in
% the data array, and then calling the generated nested setter for the payload.


% All getters and setters are directly accessing the top-level boxed type.
% In the example shown in Figure~\ref{fig:cogent-type-ex}, the getters for
% its child field $\FieldN{struct}$, grandchild fields $\FieldN{a}$ and $\FieldN{b}$ will all get from $\ConsN{Example}$, since the data structure is already
% flat.
\noindent As can be seen, a getter function (or similarly a setter) incurs a data format conversion:
It turns the low-level data format into a high-level typed value in C. Therefore
getting and passing around a large unboxed type can be expensive at run-time.
In practice though, if the program is carefully designed and implemented,
and the programmer only passes the minimal structures needed to external
functions, there is a good chance that an optimising C compiler is able to
eliminate unnecessary data conversions. \cogent allows programmers to
annotate functions with a \texttt{CINLINE} pragma, so that the \cc|inline|
modifier is generated in the C code, exposing more optimisation
opportunities. 
%In future, we plan to extend \cogent's uniqueness type system
%to allow for pointer access to unboxed types, improving the run-time 
%performance and reducing reliance on the C optimiser in these cases.


% As mentioned above, nested field getters and setters are generated 
% when nested (unboxed) records are involved, like \lstinline|struct| in Figure~\ref{fig:cogent-type-ex}.

\subsubsection{Invalid Bit Patterns}
Calling a getter on an external word array can lead to unexpected behaviours
when the
data format is invalid. This can happen when variant types are involved, if
the value in the tag part does not match any tag values declared in the
\dargent layout. Any undefined tag values will be treated as the last
constructor's tag. For example, if the layout of a variant
is $\DLVar{2\bit}{\ConsN{A}(0) : 1\byte, \ConsN{B}(1) : 2\byte,
\ConsN{C}(2) : 4\byte}$
and the tag value seen in the data format is 3, which does not match
any defined tag values but also fits in
the 2-bit space reserved for the tag, it will be deemed as the last
alternative, namely $\ConsN{C}$.

\dargent is not a low-level data parsing language, nor an interface language,
therefore the generated C getters do not check for validity, and
assume that the input is valid. Users of the language are responsible for
checking the validity of any incoming data.

\subsubsection{Required Properties of Getters and Setters}
As mentioned above, the compilation of a \cogent record type $T$ to C
can be arbitrary as long as getters and setters are provided. When it
comes to formally verifying the compiled C program, these getters and setters
are expected to compose well:
\begin{enumerate}\item 
  setting a field does not affect the result of getting another field;
\item getting a set field should return the set value (or at least an
  equivalent value).
\end{enumerate}
%
These properties are explained 
in~\autoref{sec:get-set-lemmas}. It is worth mentioning that, somewhat unintuitively,
extending the verification
framework to account for \Dargent does not require proving that
the generated getters and setters are accessing data at the locations
specified by the layout.

\subsection{Implementation}

As mentioned in~\autoref{ss:dargent-core}, the \dargent core calculus initially represents a bit range
as a pair of integers denoted by $\bitrange{o}{s}$, where $o$ indicates the offset (in bits)
to the beginning of the top-level datatype in which it is contained, and $s$ indicates how many
bits it occupies. This representation is concise and easy to work with when typechecking the core language.
Such a representation, however, does not necessarily lend itself to an easy code
generation algorithm. Therefore we have another
step to convert each bit range into a list of \emph{aligned bit ranges} tailored for C code generation. Each aligned bit range is essentially a word.

An aligned bit range $\abr{w}{o}{s}$ is a triple of integers, where $w$ is the
word-offset to the top-level datatype, $o$ is the bit-offset inside this word,
and $s$ is the number of bits occupied.
Each aligned bit range contains the information of how the bits in each word is used.
The generated C getter/setter consists of a series of statements, each operating on one
word via bitwise operations.
Each such C statement is generated according to one aligned bit range. This one-to-one mapping simplifies
C code generation. Because of the deployment of aligned bit ranges,
it is easy for us to choose an appropriate word size (namely the size of the
aligned bit ranges) specifically tailored for the application in question.
This flexibility is important for low-level programming (see \autoref{ssec:power-control-system} for a concrete example) and for high performance.
% The intuition is that machine instructions typically work on the granularity of machine words.
% In the generated C code, we want to mimic that.
% have to partition the layout encoded in a bit range into words, each represented as
% an aligned bit range, so that they closely correspond to the C statements we are going to generate.
% ~\footnote{This is with the exception of array compilation, which we will cover shortly.

% For this reason, we break down a bit range into smaller units, thus each one can be operated on by
% a single statement in C.
% The granularity, we call \emph{aligned size}, is thus set to be a machine word.
% (but in general can be different, as we shall see when coming to arrays).
% One key meta-property that holds on aligned bit ranges is as follows:
% $$\forall \abr{w}{o}{s}.\, w \ge 0 \wedge o \ge 0 \wedge s \ge 1 \wedge w + o \le A$$
% where $A$ is the align size.

% The compilation of arrays is handled slightly differently. In the current language
% design, we restrict the layout of each element to be byte-aligned. It means that
% each element will have to be at least one byte long. We thus set the granularity to
% a byte, rather than a machine word, when we encounter arrays.



\section{The Verification Framework with \Dargent}
{\label{sec:verification}}

%As we saw in~\S\ref{ss:formalised-semantics}, 
As we have demonstrated, \Dargent affects both
the formalisation of the type system and the C code generation of the compiler.
Moreover, it also affects the generated correspondence proof between
the C code and the pure shallow embedding in Isabelle/HOL. For this proof,
only the refinement between \Cogent's stateful update semantics
and the generated C code must change.
Above this low-level layer, the functional embeddings of \cogent still use
high-level algebraic types, regardless of the specified layouts.
%by the way data
% In addition to the extension of the type system, which impacts
% the deep embedding,
% the main component that must be adapted is the proof of
% correspondence between the update semantics and the C code. The crucial
% change is that getting or setting a field in \cogent does
% not correspond anymore to getting or setting a corresponding C field,
% but to calling a custom getter or setter.
%Indeed, at the C level, a field is now encoded in an array of words
%according to the layout.  Getting or setting a field in \cogent corresponds to
%a call to a generated custom getter or setter, in C.
In the following subsections,
we discuss  the extensions to the verification framework regarding:
%as \cogent records with layouts now compile to flat arrays:
%In the following subsections, we discuss what changes to the verification
%framework were necessary in this correspondence statement to take \Dargent into account:
\begin{enumerate}
  \item the correspondence between \cogent record values and C flat arrays;
  \item the correspondence between record operations in \cogent and C;
  \item the compositional properties of generated getters and setters;
  \item the correspondence between generated getters/setters and the specified
     layout.
\end{enumerate}
These extensions to the verification framework not only ensure that the C code
generated by the compiler is a refinement of the \cogent code, but also
prove that \cogent types are laid out in C correctly according to the programmer's \dargent
specification. Again, all the results we present in this section are
developed and checked in Isabelle/HOL.


The verification of any software system must assume the correctness of a 
\emph{trusted computing base}. In \cogent systems, this trusted computing base is largely
comprised of externally provided C code~\citep{OConnor:2016}. While the
\cogent framework allows for manual verification of this C 
code~\citep{Cheung_OR_22}, with \dargent we eliminate large swathes of this C code entirely, specifically the so-called ``glue code'' which translates between data formats. This reduces the size of the TCB without imposing any additional verification burden.
 


% In a last subsection, we describe additional generated proofs ensuring that C field getters and
% setters respect the given record layout specification.

% \review{what happens? what are the
%   implications? why not tackled? $\times$2.}

% \review{what proofs are currently generated, what proofs are still manual,
% what is still trusted not at all proven, and what will be done versus what will be left trusted. what guarantees Dargent currently provides and what could go wrong. which lemmas stated in the paper are formally proven. when something is not yet supported, what currently happens and the implications of that.}

% \review{are the lemmas formally proved? }

\subsection{Relating Record Values}\label{ssec:record-value-rel}

The correctness of compilation is justified by a \emph{refinement proof} that states that a simulation relation, called a \emph{refinement relation}, holds between the source language and the target language.
The refinement relation from the \cogent update semantics 
and the compiled C program is defined in terms of a relation between
\cogent values (in the sense of the update semantics)
and C values. 
Without \Dargent, this value
relation is straightforward for records: A \cogent record relates to a C structure with the same fields, such that each \cogent field is itself value-related to 
the corresponding C field.

With layouts, a \cogent record now relates to (essentially) a flat word array, and each
\cogent field relates to the result of applying the relevant custom
getter to the array.
This requires an appropriate embedding of our getters in Isabelle.
To this end, we have implemented an Isabelle/ML procedure that imports the C getters and setters generated for each field using the AutoCorres library, which produces a monadic, shallow HOL embedding that models the semantics of the C functions. Our ML procedure then simplifies these embeddings to pure Isabelle functions.
We call these simplified functions
\emph{direct getters} and \emph{direct setters}.
These direct getters and setters allow a simple statement
of the compositional properties between getters and setters,
detailed in~\autoref{sec:get-set-lemmas}.
%\subsubsection{Overview of the value relation}

% The correspondence statement between the deep embedding with update
% semantics and the AutoCorres embedding~\citep[Definition 2]{OConnor:2016} 
% relies on the definition of a value relation,
% for each C type, between:
% \begin{itemize}
%   \item a value, in the sense of the update semantics, and
%   \item an Isabelle term whose type represents the C type according to
%     AutoCorres.
% \end{itemize}
% When embedding C code in Isabelle using the AutoCorres library,
% each C type is represented directly as an Isabelle type. For example, the C structure
% \lstinline[language=C]${ U8 a; U16 b; }$ essentially embeds as
% an Isabelle record type consistings of two fields $a$ and $b$ with the types 
% \lstinline$8 word$ and \lstinline$16 word$ of natural numbers fitting on
% 8 and 16 bits respectively.
%  Then, C functions are embedded as \emph{monadic} Isabelle programs for the monad
%  \texttt{NonDet}, allowing for failure, non determinism, and stateful operations. 
% For example, a C
% function \texttt{U16\ square(U8\ a)} is embedded as a term of type
% \texttt{8 word $\Rightarrow$\ (16 word) NonDet} in Isabelle. A main motivation
% of \cogent is actually to get a pure non-monadic embedding.

% The value relation is defined automatically (by some ML procedure) for
% each involved C type in the generated C code. When no custom layout is
% involved, the value relation for a C struct between a value
% \texttt{v} and a term \texttt{t} of type \texttt{T} (representing the
% C type) enforces \texttt{v} to be a record with the same number of
% fields as \texttt{T}, such that the field values are themselves
% value-related to the corresponding fields of \texttt{t}. The adequacy of
% this generated value relation relies on the fact that records are
% straightforwardly compiled to C structures, when no layout is involved.

% \subsubsection{Records with custom layouts}

% When assigned a layout $\ell$, a \cogent record type
% %\texttt{T =  \{\ a\ :\ A,\ b\ :\ B\ \}}
% \lstinline$T = { a : A, b : B }$ translates (essentially) to an array of
% words \texttt{$\sem{ T }_\ell$}. The corresponding value relation between
% a value (in the sense of the update semantics) and a
%  word array \texttt{arr} still enforces the value to be a record type with the same number of fields as
% in \texttt{T}, and moreover that the field values relate
%  to the corresponding data encoded in the array
% of words according to the layout.

% In the particular example considered here, the value relation would enforce that the first
% field of the value relates to \texttt{get\_a\ arr}, where
% \texttt{get\_a} is the custom getter whose expected type is
% \texttt{$\sem{ T }_\ell$\ $\Rightarrow$ $\sem{ A }$}.

% \subsubsection{Direct definition of custom
% getters}
% \label{ss:direct-get}
% AutoCorres provides an Isabelle definition of the custom getters from
% the C code, but they do not have the right type for the value relation, as
% they are monadic. In the example above, it provides
% \texttt{get\_a\ :\ $\sem{ T }_\ell$\ $\Rightarrow$ $\sem{ A
%   }$\ NonDet}. Defining the value relation as described above thus requires 
% a \emph{direct} definition \texttt{direct\_get\_a\ :\ $\sem{ T
%   }_\ell$\ $\Rightarrow$ $\sem{ A }$} to be inferred. This is done by a ML procedure which
% inspects the monadic definition and simplifies it.

%There are other possible ways to achieve this goal. For example, the
%compiler could be extended to generate these direct getters in Isabelle.


\subsection{Relating Record Operations}\label{sec:record-lemmas}

Without \Dargent, records directly compile to C structures, so
it is trivial to relate
\cogent record operations to their compiled C versions
and show that they correspond.
With layouts, the correspondence proof becomes slightly more involved, since it relates
getting or setting fields in \cogent to calling custom getters or setters
in C. These correspondence statements are proved automatically by custom tactics we have developed 
which exploit the compositional properties of getters and setters, detailed in the next section,
as well as the relation between their monadic and direct embeddings.


\subsection{Verifying the Custom Getters and Setters}
\label{sec:get-set-lemmas}

To relate \cogent record operations and operations on a C array
$\texttt{arr}$ that models the record type, getters and setters are
expected to obey the following compositional laws schematically
summarised below. Given an $\texttt{arr}$ array that models a record
type:
\begin{description}[leftmargin=\parindent,labelindent=\parindent]
  \item[\text{[Roundtrip]}]
    If a \cogent value $x$ relates to a C value $v$,
    then it must also relate to the C value
    $\texttt{get\_a\ (set\_a\ arr\ $v$)}$.
    This is a weakening of the more intuitive roundtrip statement
\texttt{get\_a\ (set\_a\ arr $v$)\ =\ $v$}, which does not always hold:
As a counterexample, consider a boolean field. Since C does not natively provide such a 1-bit type,
booleans are typically represented using a single byte \lstinline{U8}. But a custom getter
for a boolean field would always return 0 or 1, picking 
the relevant bit as specified in the layout. Thus, we obtain a counter-example by taking $v=3$.
  \item[\text{[Frame]}]
    For distinct fields \texttt{a} and \texttt{b},
    \texttt{get\_a\ (set\_b\ arr\ v)} is equal to \texttt{get\_a\ arr}.
\end{description}
These statements are subject to typing constraints that we do not detail for brevity.
%Here, \texttt{direct\_get/set} refers to the direct definition of
%getters/setters mentioned in~\S\ref{ss:direct-get}, while \texttt{ac\_get/set} refers to
% the monadic AutoCorres embedding of them.
% Direct setters are inferred by a similar process as the direct getters.
% \lstinline${ sum : < A U8 | B U8 > }$ with a layout assigning the same location
% to the payloads of both alternatives. The custom setter takes as arguments a word array 
% and a structure with three fields: one for the tag value, one for the A-payload,
% and one for the B-payload (this is indeed how \cogent variant types are compiled). 
% It puts one payload or the other, depending on the tag value, at the 
% location specified by the layout. This is clearly not injective in the second argument, and thus
% there is no chance to get the intuitive strict equality.

% Another example where the equality does not hold is that of bit fields, involving integers with a non 
% standard width, for example $\ConsN{U29}$. 
% Such an integer type does not exist natively in C (neither in \cogent, see~\S\ref{ssec:bitfields}
% for an account of them using abstract types), but it can be represented
% as a $\ConsN{U32}$. However, getting and setting a field of this type with a 
% $\ConsN{U32}$ value results in its truncation on twenty-nine bits.

\subsection{Ensuring Getters and Setters Comply with Layouts}
\label{ss:sanity-checks}

As noted before, our compiler correctness certificate
does not require that our generated getters and setters actually respect the layout specification---only that they satisfy the compositional properties stated previously.
Thus, it remains to ensure that the generated getters and setters comply with data layouts specified by the user.

% Writing a generic specification of getters and setters
%once and for all, that does not need to be generated for each program (and carefully
%checked by the user), requires a deep embedding of them.
% AutoCorres relies under the hood on the C parser library that could be used for this purpose.
% However, we adopt another strategy:
To this end, we  generate specifications for getters in the deep embedding of \cogent,
returning \cogent values (in the update semantics) from a word array,
according to the layout.
More specifically, we implemented a
generic getter \lstinline|uval_from_array|
as a recursive function in Isabelle that takes as input a \cogent type, a
\Dargent layout,
a word array\footnote{Each of these inputs is of a suitable type inductively
  defined in Isabelle.},
and retrieves the value of the field of the given
\cogent type from the word array according to the given field layout.



We designed tactics to automatically show that generated getters respect this
specification. More precisely, given a field $b$ of \cogent type $\tau$ with associated layout
$\ell$, the tactics
prove that for any word array $a$ that has a valid
format according to the layout, \lstinline|get_b|~$a$
relates to \lstinline|uval_from_array|~$\tau\ \ell\ a$ in the sense of the value
relation,
 where \lstinline|get_b| is the 
direct getter.

We have not yet done the same for setters, although the compositional properties we verify,
together with the correctness of getters, already provide some formal guarantees. Additionally, we also automatically prove explicitly that
setters do not flip any bit that is outside the field layout. 

%\subsection{Tactics for the get/set lemmas}
%We wrote one ML tactic for each kind of get/set lemmas. They are about 50 lines long 
%each and rely on the tactics \texttt{word\_bitwise} and \texttt{monad\_eq} provided by the
%AutoCorres library, that simplify equalities between integers and monadic programs,
%respectively.
% Optimising them is left for future work.
% The first get/set lemma mentioned above
% can be unnecessarily slow to prove in the case where multiple variant types are involved,
% as the tactic currently checks every possible combination of constructors.

\subsection{Endianness}\label{ssec:endianness-verif}

Since \Cogent relies on the AutoCorres library, which does
not support big-endian architectures yet, we assume that the machine 
is little-endian.
Therefore, only big-endian annotations need to be seriously accounted for: 
Big-endian fields of type $\PrimType{U16}$, $\PrimType{U32}$, and $\PrimType{U64}$ require reversing the bytes on retrieval and setting. This reversal is performed in dedicated optimised functions.
As an example, here is the code for reversing the bytes in a $\PrimType{U32}$:
\begin{lstlisting}[language=C]
  static inline u32 swap_u32(u32 v) {
    v = ((v << 8) & 0xFF00FF00) | ((v >> 8) & 0xFF00FF);
    return (v << 16) | (v >> 16);
  }
\end{lstlisting}
This additional reversal step does not require any
extension to the proof tactics for the compositional properties of getters and
setters described in \autoref{sec:get-set-lemmas}, as they proceed first by unfolding all involved function definitions and applying an automated simplifier for bitwise operations, which can automatically prove that these swap operations are an involution. 
Recall that the value relation described in \autoref{ssec:record-value-rel} 
is defined based on our generated custom getters, which already take endianness into account, so the value relation also naturally takes 
endianness into account. 

 The proof described in \autoref{ss:sanity-checks}, that our getters respect the given layout, also must accommodate endianness. 
Specifically, we use a purely functional HOL embedding of the swap function to 
specify getters of big-endian values. Because our HOL embedding follows the
C implementation closely, our proof tactics also need no extension for big-endian fields to prove correctness of getters.

We also prove as a sanity check that these swap functions are correct, in the sense 
 that they reverse the byte order.


\section{Applications}
\label{sec:applications}

\review{None of the examples use variants, which is arguably the most
important feature that \dargent provides which surpasses the expressiveness
of C.}

\ambroise{we do use variants (to account for enums), although with empty constructor payloads. If I
  have time to implement it, the
  maybe example suggested by the reviewer could be more convincing.
}

In this section, we showcase how \Dargent helps in developing and formally
verifying systems code via some small examples, the full source code of which is available in the supplementary material.
 
% \review{more information about effort involved (especially something quantitative, if possible), and possibly information about performance. also for \S~6.}


\subsection{Power Control System}\label{ssec:power-control-system}

To demonstrate the improved readability of programs that \Dargent offers, 
we reimplemented the power control 
system for the STM32G4 series of ARM Cortex
microcontrollers by ST Microelectronics, based on a C implementation
from LibOpenCM3\footnote{\url{http://libopencm3.org/}}, an open-source low-level
hardware library for ARM Cortex-M3 microcontrollers.
The original C file\footnote{\url{http://libopencm3.org/docs/latest/stm32g4/html/pwr_8c_source.html}}
is about one hundred lines long, and consists of eight functions of a few lines
each that perform some bitwise operations on the device register. 
% \review{how long is the corresponding \cogent code?}
The \cogent implementation is a bit smaller as each function is now one line long. The
most involved part is the \dargent layout itself, which is derived from the hardware specifications of the device.

The 32-bit power control register holds several pieces of information, each occupying a certain
number of bits in the register. We define the register as a record type in \cogent
and use \dargent to prescribe the placement of each field. With a \dargent layout,
the functions in the power control code no longer involve bitwise operations, but merely set values to individual fields.
For example, the function to set the voltage output scaling bit in C:
\begin{lstlisting}[language=C]
  void pwr_set_vos_scale(enum pwr_vos_scale scale) {
    uint32_t reg32;
  
    reg32 = PWR_CR1 & ~(PWR_CR1_VOS_MASK << PWR_CR1_VOS_SHIFT);
    reg32 |= (scale & PWR_CR1_VOS_MASK) << PWR_CR1_VOS_SHIFT;
    PWR_CR1 = reg32;
  }
\end{lstlisting}
translates to the \cogent function setting the $\FieldN{vos}$ field in the record:
\begin{lstlisting}[language=Cogent]
  pwr_set_vos_scale : (Cr1, Vos_scale) -> Cr1
  pwr_set_vos_scale (reg, scale) = reg { vos = scale } 
\end{lstlisting}
It can be seen from this example that \cogent allows systems programmers to write code
on an abstract level, while retaining low-level control of the implementation details with \dargent{}.


% for the STM32G4 series of ARM Cortex Microcontrollers by ST Microelectronics.

% \review{Furthermore, while Dargent provides a lot of control over data *layout*, it does
% not (seem to) provide control over the exact *access pattern* that is used to
% get and set fields. Of course, usually that is an implementation detail, but
% several of the examples involve using Dargent to describe the layout of hardware
% registers (which I assume refers to hardware that is controlled via MMIO).
% Hardware register access in C is usually done with `volatile`, ensuring that the
% access happens exactly as written (e.g., a single 32-bit write, even if only the
% first 16 bits were changed). For example: presumably, in pwr\_set\_vos\_scale,
% PWR\_CR1 is a volatile 32-bit integer? Does the code generated by Dargent
% guarantee that the setter will do a single 32-bit volatile write? The
% description of the generated setters and getters is rather vague, so it might
% well be the case that Dargent generates a 16-bit write when only some bit in the
% first half of PWR\_CR1 changed, which might or might not actually be supported be
% the hardware this driver is controlling. Or, worse, it could generate multiple
% small writes. In these settings where the exact sequence of memory writes and
% their sizes matters, how does a Dargent user ensure that the generated code does
% what it is supposed to do?}

% \zilinc{Do we talk about it here, or early in \S~3?}
% 
% \ambroise{I think we should sketch the access pattern in sec 3.4 (and discuss
%   performance, in particular regarded accessing nested subfields) on custom
%   getters and setters, but discuss the volatile issue here.}


%\paragraph{Dealing with volatile registers}
%This example also demonstrates a current limitation of our system:
%the verification framework does
%not provide any guarantees about the exact access pattern of
%getters and setters. 
%This is however crucial when dealing with hardware registers, since
%the write (or read) order can matter.
% Consider, for example, the following
%\cogent function from the power control system.
%\begin{lstlisting}[language=Cogent]
%pwr_enable_power_voltage_detect : (Cr2, Pvd_level) -> Cr2
%pwr_enable_power_voltage_detect (reg, pls) = reg { pls = pls, pvde = True }
%\end{lstlisting}
%This code actually compiles to two subsequent writes to the register CR2:
%it first sets the programmable voltage detector level \lstinline|pls|
%and then enables the programmable voltage detector (by setting the boolean
%flag).
%If the writes were written in reverse order, however, the detector would run
%with an unknown selection level between the two writes.
%% Although the device specification does not explicitly
%% discourage such a write pattern,
%To avoid this kind of unwanted behaviour, sequentialisation is
%usually avoided when dealing with hardware registers. 
%In fact, the original C implementation performs a single write to the register:
%\begin{lstlisting}[language=C]
%void pwr_enable_power_voltage_detect(uint32_t pvd_level)
%{
%  uint32_t reg32;
%  reg32 = PWR_CR2;
%  reg32 &= ~(PWR_CR2_PLS_MASK << PWR_CR2_PLS_SHIFT);
%  PWR_CR2 = (reg32 | (pvd_level << PWR_CR2_PLS_SHIFT) | PWR_CR2_PVDE);
%}
%\end{lstlisting}
%%In future work, we plan to extend \cogent so that it can perform
%%multiple field updates/reads in a single operation, when possible.

% The difficulty
% Beyond the implementation, a  difficulty relates to
% the formalisation of the different write patterns, which are
% currently indiscernable because of the memory model of the
% AutoCorres library on which \cogent's verification framework
% is based.

% As an example, flipping multiple booleans flags
% will always result in multiple writes
% However, hardware register access typically requires 
% that  depends C is usually done with `volatile, ensuring that the
% access happens exactly as written (e.g., a single 32-bit write, even if only the
% first 16 bits were changed).



\subsection{Bit Fields}\label{ssec:bitfields}

In systems code, there is a regular pattern consisting of declaring a structure whose fields have non-standard
 bit widths, as in the following example taken from a CAN driver%
 \footnote{\url{https://github.com/seL4/camkes-vm-examples/blob/89f5d7b7ac373c8e9f000e80b91611e561358ef6/apps/Arm/odroid_vm/include/can_inf.h\#L34}}
 where the first field is an identifier on 29 bits.
\begin{lstlisting}[language=C]
  struct can_id {
    uint32_t id:29;
    uint32_t exide:1;
    uint32_t rtr:1;
    uint32_t err:1;
  };
\end{lstlisting}
Although this is not supported by the AutoCorres library~\citep{autocorres} on which the \cogent
verification framework is based, we can simulate this feature using \Dargent.

% As a first attempt, consider the following \cogent type.
% \begin{lstlisting}
%   type R = { f1 : Bool, f2 : U8 }
%      layout record { f1 : 1b, f2 : 4b }
% \end{lstlisting}
% The \cogent compiler would actually complain that a \lstinline{U8} value does not fit in 4 bits.
% % As far as the verification is concerned,
% This mismatch compromises the first get/set
% lemma (see~\S\ref{sec:get-set-lemmas}), since setting the field \lstinline{f2}
% with a \lstinline{U8} value $v$ and then retrieving
% it does not yield $v$, but its truncation on 4 bits.
% To solve this issue, we need to introduce a new type \lstinline{U4} in \cogent. 
% This type essentially translates to \lstinline{U8} in C, but the value relation 
% for this type ensures that the value fits on 4 bits.
The last three fields are one bit long and can thus be handled with a \cogent
boolean type.
In order to cover the 29-bit field, we can use a primitive 29-bit integer $\PrimType{U29}$. These non-word-size integers are a new extension we added to \cogent for this purpose, and consisted only of a dozen changed lines of code. We compile such integers to
the smallest standard integer type that can contain the type, since the C language does not
natively support such types. For instance, $\PrimType{U7}$ is
compiled to $\PrimType{U8}$, while $\PrimType{U20}$ is compiled to $\PrimType{U32}$.
Thanks to the extension of the value relation (see \autoref{ssec:record-value-rel}) 
to those new types, compiled \cogent programs maintain the invariant that C
values always fit in the narrower bit-width. 
%
%
%
% However there is no built-in type
% that can be used for the 29-bit field (recall the match rules in~\S\ref{ss:formalised-semantics}). 
% Thankfully, \cogent provides a modular mechanism to introduce new types, as abstract types,
% whose definitions must be provided by the programmer in C.
% The verification framework must also be extended to take into account abstract values for these types.
% % , and provide the value relations for these new types. 
% For example, in the shallow embedding, we embed $\ConsN{U29}$
% as the Isabelle type of natural numbers that fit in 29 bits.
%
% By first declaring the abstract type $\ConsN{U29}$,
With a \dargent layout, we can define a \cogent version of the above C structure:
\begin{lstlisting}[language=Cogent]
  type CanId = { id : U29, exide : Bool, rtr : Bool, err : Bool }
    layout record { id : 29b, exide : 1b, rtr : 1b, err : 1b}
\end{lstlisting}
% Abstract types are implicitly boxed, hence the additional
% unboxed annotation \lstinline{#} in front of $\ConsN{U29}$.

% When encountering an unboxed abstract type in a layout, the compiler does not check
% any compatibility between the bit range and the type. In the generation of getters and setters,
% the field is dealt as if it were
% an ordinary integer type, with a final C cast to the abstract type.
% For some technical reason, abstracts types must always be embedded in
% a custom structure. As a consequence, we introduce an additional layer of compilation 
% that replaces this brutal C cast with a more explicit one.
% Nicely enough, our tactics are still able to prove the get/set lemmas
% stated in~\S\ref{sec:get-set-lemmas}.
% 

% To demonstrate how this approach works in practice, we proved functional correctness
% of the following simple function that updates the second or third field
% depending on the value of the first field.
% \begin{lstlisting}
%   foo : R -> R
%   foo r =
%     if r.f1 !r then
%       let v = u8_to_u4 (u4_to_u8 r.f3 .&. 0x0c) !r in
%       r {f3 = v}
%     else
%       let v = u8_to_u2 (u2_to_u8 r.f2 + 1) !r in
%       r {f2 = v}
% \end{lstlisting}
% Abstract types alone are not very useful since no \cogent operation applies to them.
% We thus extend our example with back and forth casts between $\ConsN{U29}$, and $\ConsN{U32}$,
% declared in \cogent as \emph{abstract functions}. This means that we implement them in C, 
% and we also provide their semantics at each layer of the refinement proof, and 
% prove properties about them~\citep[\S3.3.2]{OConnor:2016}.
% \begin{lstlisting}
%   u4_to_u8 : #U4 ->  U8
%   u8_to_u4 :  U8 -> #U4
%   u2_to_u8 : #U2 ->  U8
%   u8_to_u2 :  U8 -> #U2
% \end{lstlisting}
% These functions are abstract: they are implemented in C. We exploit the design
% of the verification framework that eases the extension of the semantics to account
% for abstract functions. 

% \review{how much manual effort is needed? hours? LoC? etc.}
% Although some manual verification effort is required (about 700 lines of Isabelle/HOL proof) on these irregular-sized integers and the
% cast functions, it is only a one-off effort, and they can be mechanically produced.
% We leave it as future work to have native support for these integer types from the \cogent language
% and the compiler.


\noindent Bit fields can play an important role in systems programming. This example demonstrates that 
\dargent not only allows programmers to store data in a compact manner, but also provides a 
principled way to reason about C bit fields without needing to extend our C verification tools
to support the bit field feature in C.


\subsection{Custom Getters and Setters}\label{ssec:partial-rec}

Suppose we are writing a \cogent system that must work on a complicated, externally defined C structure that involves many fields, but the \cogent component is only concerned with a few specific fields. 

Because \dargent makes all record accesses go through getter and setter functions, we can reuse this mechanism by providing \emph{custom} getter and setter functions in C, that extract (or write to) the relevant fields inside the large C structure.  Then, we can represent this large C structure as a simple \cogent record with only the relevant fields.
%% \begin{lstlisting}[language=Cogent]
%%   type Entry = { id : U32, value : U32 }  
%% \end{lstlisting}
The custom getters and setters provided by the user must compose well (see~\autoref{sec:get-set-lemmas}).

We have successfully applied this approach on a small example, where the
C structure has the same fields as the original \cogent type, extended
with an additional field. Whilst this example is contrived, this feature
does have real application: It makes it possible to \emph{inherit}
pre-defined C data structures. For example, in a previous \cogent implementation of a Linux \texttt{ext2fs} driver, \citet{Amani:2016}
had to define their own \Cogent \verb!Ext2Inode! type which corresponds to the standard C \lstinline[language=C]!struct ext2_inode! type, but did not include fields that were irrelevant to the \cogent implementation, such as  certain unsupported file modes and spinlocks. This required tedious glue code to marshal back-and-forth between the representations. 
Using this technique, the standard C \lstinline[language=C]!struct ext2_inode! can be used directly as the representation of the \verb!Ext2Inode! type, thus eliminating this glue code entirely.


%Currently, this customisation of getters and setters is accomplished rather awkwardly, requiring the programmer to assign the type a dummy layout to force the generation of getters and setters, and then replace the generated getters and setters with the desired custom functions.
%Of course, to get the full refinement proof, one needs to prove that getters and
%setters compose well (\S~\ref{sec:get-set-lemmas}).
% In future
%work, we plan to improve the compiler to make this approach more
%user-friendly.

\section{Verification of a Timer Device Driver}
\label{sec:timerdriver}

In this section, we present a formally verified \cogent timer driver (available
in the supplementary material), that we ran
successfully on an ODROID hardware, based on the seL4 operating system~\citep{seL4:2022}. The formal verification was conducted in Isabelle/HOL.
Our formalisation took advantage of the fact that the
shallow embedding of the \cogent program in Isabelle remains simple, as the layouts are fully transparent to the functional
semantics of the \cogent program.


Our \cogent implementation is based on a C driver\footnote{\url{\urltimerdriver}}
% that
% is sixty lines long and 
consisting of an interface for two timers, called A and E, provided by the device.
Both implementations are about the same size ($\approxeq$60 LoC, excluding type
and layout declarations). 
%
The driver can be used to:
\begin{itemize}
\item 
  measure elapsed time since its initialisation, using the device timer E, and
  \item
    generate an interrupt at the end of a (possibly periodic) countdown,
    using the device timer~A.
\end{itemize}
The driver consists of four functions:
\begin{itemize}
\item
  \href{\urltimerdriver#L13}{\lstinline{meson_init}}
  initialises the device register;
  \item
    \href{\urltimerdriver#L29}{\lstinline{meson_get_time}} returns the elapsed time
    since the initialisation (in nanoseconds);
  \item 
    \href{\urltimerdriver#L29}{\lstinline{meson_set_timeout}}
    sets the countdown value, making it possibly periodic;
  \item 
    \href{\urltimerdriver#L59}{\lstinline{meson_stop_timer}} stops the countdown.
\end{itemize}
The driver state is passed around as a C structure called
\href{\urltimerdriverh#L82}{\lstinline{meson_timer_t}}, consisting of
the location of the device registers in memory and a boolean flag
remembering whether the countdown is~enabled.


\subsection{Using \dargent}
In the original C implementation, operations on the timer registers are largely based on
bitwise operations, which are unintuitive to read, error-prone, and more difficult to prove correct. Modelling the timer registers as algebraic
data types, like records and sum types,
makes the program easier to read and to reason about, while \dargent still allows us control over low-level representation details. 
%
With \dargent, the timer register can be defined as:
\begin{lstlisting}[language=Cogent]
  type Meson_timer_reg = {
     timer_a_en : Bool,
     timer_a : U32,
     timer_a_mode : Bool,
     timer_a_input_clk : Timeout_timebase,
     timer_e : U32,
     timer_e_hi : U32,
     timer_e_input_clk : Timestamp_timebase
  } layout record {
       timer_a_mode : 1b at 11b,
       timer_a_en : 1b at 15b,
       timer_a_input_clk : LTimeout_timebase,
       timer_e_input_clk : LTimestamp_timebase,
       timer_a : 4B at 4B,
       timer_e : 4B at 72B,
       timer_e_hi : 4B at 76B
    }
\end{lstlisting}
Contrary to the C definition, which is mostly type-less and relies on macros and bit-shifting
to denote the semantics of the bits in use, the \cogent definition comes with a lot more
type information, and also concisely prescribes how bits are used and placed.
A typical set-value operation in C can thus be rewritten in a more abstract way in \cogent as
record field updates, shown in~\autoref{fig:timer-set}.


 % Not only is it more readable,
 % but it also induces a nicer Isabelle shallow
 % embedding, using records and variant types, easier to formally verify.

%   As an example of layout specification, the following one
%   details the memory representation of the variant type for the timer A units:
%   \begin{lstlisting}
%     layout LTimeout_timebase =
%     variant (2b) {
%         TIMEOUT_TIMEBASE_1_US  (0) : 0b
%       , TIMEOUT_TIMEBASE_10_US (1) : 0b
%       , TIMEOUT_TIMEBASE_100_US(2) : 0b
%       , TIMEOUT_TIMEBASE_1_MS  (3) : 0b }
%   \end{lstlisting}
%   Except for the shared tag location (2 bits long), each payload
%   is 0 bit long since the constructors do not take any argument. 
%   This layout corresponds to the following C defined constants:
%   %\footnote{\url{\urltimerdriverh}}:
%   % in the interface  file:
%   \begin{lstlisting}[language=C]
%     #define TIMEOUT_TIMEBASE_1_US   0b00
%     #define TIMEOUT_TIMEBASE_10_US  0b01
%     #define TIMEOUT_TIMEBASE_100_US 0b10
%     #define TIMEOUT_TIMEBASE_1_MS   0b11
%   \end{lstlisting}
%   
  
\subsection{Verification}
To formally verify the driver, we first implemented a purely functional
specification of the driver. Then, we proved that the shallow embedding
of the \cogent driver refines it.

Both the device and the driver states are threaded through the functional implementation.
For example, the initialisation function performs a functional update of the device state,
represented as a record similar to
\lstinline|Meson_timer_reg| (ignoring the layout annotations):
\begin{lstlisting}[language=isa]
  initialize s =
    (s \<lparr> deviceState := (deviceState s)
          \<lparr> timer_e_low_hi := 0,
             timer_a_input_clk := TIMEOUT_TIMEBASE_1_MS,
             timer_e_input_clk := TIMESTAMP_TIMEBASE_1_US
           \<rparr> \<rparr> )
\end{lstlisting}

% \review{numbers on effort on manual verific\ation?}
Both the specification and the manual functional correctness proof are $\approxeq$150 lines each.
The manual proof is established by straightforward equational reasoning---much easier than reasoning about the bitwise operations implemented in C. Layouts are transparent on the shallow embedding level: \cogent records are encoded as Isabelle records just as if no layouts were~specified.

This manual functional correctness proof composed with the automatically generated compiler certificate establishes
the correctness of the compiler-generated C code. All the tedium in the layout details is successfully hidden by our automatically generated compiler proofs.

\subsection{Discussion}
\label{ssec:issues}


%\zilinc{Is it relevant?}
Even though this timer driver is a short program, our formal verification nonetheless uncovered several bugs or implicit assumptions that had been made in the original C driver.

\begin{figure}
\begin{tabular}{p{0.5\textwidth}p{0.5\textwidth}}
\begin{minipage}{.5\textwidth}
\begin{lstlisting}[language=C]
  timer->regs->mux =
     TIMER_A_EN 
    | (TIMEOUT_TIMEBASE_1_MS << 
       TIMER_A_INPUT_CLK)
    | (TIMESTAMP_TIMEBASE_1_US << 
       TIMER_E_INPUT_CLK) ;
\end{lstlisting}
\end{minipage} &
\begin{minipage}{.5\textwidth}
\begin{lstlisting}[language=Cogent]
  regs {
     timer_a_en = True,
     timer_a_input_clk = TIMEOUT_TIMEBASE_1_MS,
     timer_e_input_clk = TIMESTAMP_TIMEBASE_1_US }
\end{lstlisting}
\end{minipage}
\end{tabular}
  \caption{The code 
  enables the timer A, and selects
  the time units for the timers A and E.
  \texttt{TIMEOUT\_TIMEBASE\_1\_MS} and
  \texttt{TIMESTAMP\_TIMEBASE\_1\_US} are constant macro definitions in C (left), whereas in \cogent (right)
  they can be modelled as constructors of a variant type, but compiled to 2-bit integers.}
\label{fig:timer-set}
\end{figure}

Firstly, the original implementation of the initialisation function \lstinline{meson_init}
enabled the countdown timer A, without setting a starting value for it. The behaviour of the timer
device in this case is unspecified. In fact, this countdown timer should 
only be enabled when used---that is, in the \lstinline{meson_set_timeout} function.
%We thus
%removed this instruction.
Another related issue is that the initialisation function does not ensure that
the disable flag of the driver state is synchronised with the enable flag of the
device register, but rather assumes such. We introduced a specific invariant to the functional correctness specification of this function, to make this assumption explicit.

Additionally, when verifying the function
 \lstinline{meson_get_time}, we had to explicitly assume that the timer value in the device state is not too large. While the device provides the timer value in microseconds, the function is specified to return the time in nanoseconds.
The conversion requires a multiplication by one thousand, possibly triggering
an overflow if the timer value is larger than approximately 500 years. Thus, we had to add a precondition to rule out such cases.

\subsection{Volatile Behaviour}% and pointer arithmetic}

Although some device registers (such as the configuration register)
behave as regular memory, other registers are volatile. In particular, the value of 
timer E may change between reading, so reading it twice may yield different values.
The original C implementation of \lstinline{meson_get_time} exploits this behaviour
to detect a possible overflow when reading the lower bits of timer E. 

Such volatile memory can be modelled by adapting the getter functions on those
memory locations to return non-deterministic values. While \cogent functions
are deterministic, non-determinism  can be expressed by threading an additional
abstract value through a deterministic version of the function, then
abstracting that function to a non-deterministic one without that additional
value. In particular, the getter function of a volatile register can be
replaced by a function that takes an additional abstract type as input and
returns a value of that type along with the value of the register. The shallow
embedding of this getter function can then be abstracted to a function that
ignores that additional input and returns a non-deterministic value representing
the true behaviour of the volatile register. We worked out a small example that
sums two random numbers to clearly demonstrate this idea. The verification of
this example illustrates that summing two random numbers can yield any value
(not just even values). 

%In this formalisation, we have not modelled these volatile features.
%As mentioned in Section~\ref{ssec:power-control-system}, 
%we plan to extend \cogent to handle them more carefully in future work.

% for the timer E reset: setting any value to the lower bits of timer E 
% resets the whole timer E.
% We thus implemented this reset in C, as an abstract function from the viewpoint of \cogent.
% This makes it possible to give it a custom axiomatisation in the shallow
% embedding that reflects this behaviour.
% This volatile behaviour can be modelled by introducing a layer of non-determinism in the shallow embedding, when necessary. The \cogent file system BilbyFs has an example of how non-determinism can be modelled on top of the shallow embedding~\cite{Amani:phd}. One can explore   dealing with volatile registers in a similar manner.
%In future work, we plan to address volatile behaviours in a more adapted manner, for example
%by .

%We also didn't fully re-implement original initalisation function. In particular,
%we kept on the C side the null pointer checks and the instruction finding the
%device register location, as it involves pointer arithmetic, which cannot
%be expressed in \cogent.
% \begin{lstlisting}[language=C]
% int meson_init(meson_timer_t *timer, meson_timer_config_t config)
% {
%     if (timer == NULL || config.vaddr == NULL) {
%       return EINVAL;
%     }
%     timer->regs = (void *)((uintptr_t) config.vaddr + (TIMER_BASE + TIMER_REG_START * 4 - TIMER_MAP_BASE));
%     // remaining of the initialisation function, written in Cogent
%     cogent_initialize(timer);
%     return 0;
% }
% \end{lstlisting}


\section{Related Work}\label{sec:rfw}

The idea of describing low-level data layout with high-level languages is
not new. The rich area of research makes it
challenging to fully contextualise our work within the
space. Simplifying our analysis a bit, we can roughly bifurcate the
literature into research on \emph{program synthesis} and \emph{program
abstraction}.

For instance, the PADS family of languages~\citep{FisherW11,
FisherG05, MandelbaumFWFG07}, PacketTypes~\citep{McCannC00},
Protege~\citep{WangG11}, DataScript~\citep{Back02}, Nail~\citep{Bangert_Zeldovich_14},
the generic
packet description by \citet{vanGeest_Swierstra_17}, the verified Protocol
Buffer~\citep{Ye_Delaware_19} built upon the \textsc{Narcissus}
framework~\citep{Delaware_SPYC_19}, EverParse~\citep{everparse},
and contiguity types~\citep{Slind_21} are all concerned
with synthesising a parser program (and also a pretty-printer for some of
them) from a high-level specification of the data format.

\dargent's primary focus is on the data refinement of algebraic
data types rather than securely operating on wire formats. Even though our
technology shares a lot in common, the problem we try to solve
is very different. In particular, \dargent is not a language for parsing or
converting between data formats. It is an extension to \cogent for its compiler to 
fine-tune the target code generation so that compiled code
is already in the desired format that is suitable for systems software.
In many cases, \dargent can eliminate the need for such a data marshalling tool entirely. 

Along with \dargent,
LoCal~\citep{Vollmer:2019}, SHAPES~\citep{Franco:17, Franco:19}
and \texttt{hobbit}~\citep{Diatchki_JL_05, Diatchki_Jones_06} fall in
 the program abstraction camp and 
are concerned with compiling
data structures in a program into specific layouts dictated by
the user. Programmers can therefore still work
with high-level source code, while the compiler 
does the heavy-lifting to generate the low-level mechanisms, retaining the 
separation of program logic from low-level concerns.

LoCal~\citep{Vollmer:2019} is a compiler for a first-order pure functional
language that can operate on recursive serialised data by translation into an
intermediate \emph{location} calculus, LoCal, mapping pointer indirections of
the high-level language to pointer arithmetic calculations on a base address.
The final compiler output is C code which, interestingly, preserves the
asymptotic complexity of the original recursive functions, although this
property is implementation-defined and not assured by any formal theorem.
Nonetheless, LoCal's type safety theorem does ensure a form of memory safety:
Each location is initialised and written to exactly once. The latter property
is a key difference to our work, since \Dargent can operate on mutable data by
virtue of \cogent's linear types. On the other hand, \cogent is a total
language and purposely lacks full support for recursion, we therefore do not
yet support recursive layout descriptions. Primitive recursive types for
\cogent are under development and use records~\citep{Murray:2019}. Since
we already support layout descriptions on records, we believe, once recursive
types are supported, adding support for recursive layouts would be a
straightforward engineering task. As a systems language, \cogent code often
uses abstract types such as arrays and  iteration constructs over such types.
Arrays and iteration constructs over arrays were recently verified through
\cogent's FFI~\cite{Cheung_OR_22}. We have ensured these proofs work with our
\dargent extensions. The most significant difference with LoCal is that
\Dargent is a \emph{certifying} data layout language, with generated theorems
that the translation is correct, whereas LoCal offers no verified guarantees
about its final compiler output.

The SHAPES~\citep{Franco:17, Franco:19} extension is designed for fine-tuning 
the layout of object classes for an object-oriented language to improve cache performance. 
It allows
users to define layout-unaware classes and specify what layout to use at
object instantiation time. This class parameterisation mechanism shares
some similarity with \dargent's layout polymorphism. The layouts that SHAPES
is concerned with are primarily arrays of values, which are key to better
cache locality but are not how compilers
of managed languages natively represent data in memory. Their layouts are not
down to the bit-level, but rather on the level of record fields.
In contrast, \dargent's layouts are lower-level and more flexible, and
are less tailored for a specific optimisation. SHAPES's type system maintains
memory safety properties of the program when it splits and lays out
boxed data types. Therefore it also plays a similar role as \cogent's
uniqueness type system. In our work, these two aspects are independently managed:
\dargent does not directly interfere with memory safety properties guaranteed
by \cogent's type system. 


The \texttt{hobbit} interpreter~\citep{Diatchki_JL_05} extends a Haskell-like functional language with
first-class support for bit-level types and operations (e.g.\ bit concatenation
and splitting), supporting external
representations of bit-level structures. Their work initially focused on
bit-data that can be stored within a single register and later gets extended
to memory areas realised as arrays~\citep{Diatchki_Jones_06}.
Instead of assigning a high-level type
and a low-level layout to an object in memory, their types already prescribe
the layouts, by virtue of the first-class bit-data support in the language. In
that sense, it is more comparable to the bit fields feature in the C language,
or to \cogent if the \dargent layout descriptors were subsumed by the \cogent types.
Their research novelty also lies in using advanced type system features that are readily
available in Haskell to encode the new language constructs and to perform
sophisticated typechecking, which is arguably an orthogonal
matter to data layouts.

Floorplan~\citep{Cronburg:2019} is also somewhat
relevant to \dargent, but hardly fits in either category. It is a
memory layout specification language for declaratively describing the
structure of a heap as laid out by a memory manager.
It therefore chiefly serves the implementors of memory managers rather than
systems developers and users in general, and the abstraction it provides
does not necessarily extends to algebraic types of heap objects.
The compiler follows the specification to generate memory-safe Rust code to perform
common tasks that are needed in the implementation of a memory manager.
The semantics of a heap layout
specification is denoted by the set of values that the heap can take.
In contrast, the semantics of a \dargent layout is characterised by the getter
and the setter functions.


\section{Conclusion}{\label{sec:conclusion}}

Systems code must adhere to stringent requirements on data
representation to achieve efficient, predictable performance and avoid
costly mediation at abstraction boundaries. In many cases, these
requirements result in code that is error prone and tedious to write,
ugly to read, and very difficult to verify.

We are not without hope, however, as we have demonstrated in this
paper. By using \Dargent, we can avoid the need for having the
\emph{glue} code (be it manually written or synthesised) that marshals
data from one format into another, and eliminate error-prone
bit-twiddling operations for manipulating specific bits in device
registers. Instead, we enable programmers to provide declarative
specifications of how their algebraic datatypes are laid out. Given
these specifications, our \emph{certifying} compiler generates
corresponding C code that operates on these data types directly,
 along with proofs that the generated code is
functionally correct. We have shown the applicability of \dargent on a
number of examples showcasing its support for low-level systems
features including the formal verification of a timer device driver.



\begin{acks}
Matthew Di Meglio, Sahan Fernando and Zhenyu Yao contributed to the language
design, pen-and-paper formalisation and compiler engineering during early
iterations of \dargent. We would also like to thank the anonymous reviewers
for their insightful and constructive feedback on drafts of this paper.
Robert Sison kindly offered to proofread a draft of this paper and provided many
helpful suggestions.
\end{acks}


\bibliographystyle{ACM-Reference-Format}
\bibliography{references}

\end{document}


